%
% wrspec.tex
%
% White Rabbit Spec. Currently an informal draft.
%
% best master clock algorithm 6.6.2.3
% domain number, profiles 7.1, annex J (p. 237)
% asymmetry 6.2.c, 7.4.2 (positive if t_MS > t_SM)
% 
% TODO
% Give the values of different constants in tables.
% Use thicker line for combined state in wr state machine figure.
% Take cursor out of PTP state machine figure.
\def\us{\char`\_}

\documentclass[a4paper, 12pt]{article}
%\documentclass{article}

\usepackage{fullpage}
\usepackage{pgf}
\usepackage{tikz}
\usetikzlibrary{arrows,automata,shapes}
\usepackage{multirow}
\usepackage{color}
\usepackage[latin1]{inputenc}
\usepackage{verbatim}
\usepackage{amsmath}
\usepackage{times,mathptmx}
\usepackage{chngcntr}

\graphicspath{ {../../../figures/} }

%%%%%5% used in Tomeks %%%%%%%
\usepackage{listings}
\usepackage{cancel}
%%%%%%%%%%%%%%%%%%%%%%%%%%%%%%%    fig/tomeksDrawings
%\usepackage{morefloats}

%\usepackage{amsfonts}
%\usepackage{amssymb}
%\usepackage{graphicx} 

%%%%%%%%%%%%%%%%%%%%%%%%%%%%%%%%%%%%%%%%%%%%%%%%%%%%%%%%%%%%%%%%%%%%%%%%%%%
% creating subsubsubsection notation
% src: http://www.latex-community.org/forum/viewtopic.php?f=5&t=791
%%%%%%%%%%%%%%%%%%%%%%%%%%%%%%%%%%%%%%%%%%%%%%%%%%%%%%%%%%%%%%%%%%%%%%%%%%%
\setcounter{secnumdepth}{6}
\renewcommand\theparagraph{\Alph{paragraph}}

\makeatletter
\renewcommand\paragraph{\@startsection{paragraph}{4}{\z@}%
                                     {-3.25ex\@plus -1ex \@minus -.2ex}%
                                     {0.0001pt \@plus .2ex}%
                                     {\normalfont\normalsize\bfseries}}
\renewcommand\subparagraph{\@startsection{subparagraph}{5}{\z@}%
                                     {-3.25ex\@plus -1ex \@minus -.2ex}%
                                     {0.0001pt \@plus .2ex}%
                                     {\normalfont\normalsize\bfseries}}
%\renewcommand{\thefootnote}{\fnsymbol{footnote}}
\renewcommand{\thefootnote}{\alph{footnote}}

\counterwithin{paragraph}{subsubsection}
\counterwithin{subparagraph}{paragraph}
\makeatother
%%%%%%%%%%%%%%%%%%%%%%%%%%%%%%%%%%%%%%%%%%%%%%%%%%%%%%%%%%%%%%%%%%%%%%%%%%%%%%%%

\newcommand{\eqoffset}[1]{%
  {\ensuremath{%
      {\text{offset}}_{#1}}%
  }%
}
\newcommand{\eqdelay}[1]{{\text{delay}}_{#1}}
\newcommand{\eqasymm}{{\text{asymmetry}}}

\begin{document}

\title{White Rabbit Specification: \\Draft for Comments\\\normalsize {version 2.0}\\\small{(06-07-2011)}}
\author{Emilio G. Cota\\Maciej Lipi\'{n}ski\\Tomasz W\l{}ostowski\\Erik van der Bij\\Javier Serrano}
\date{July 2011}
\maketitle
\thispagestyle{empty}

\begin{figure}[ht!]
  \centering
  \vspace{1.3cm}
  \includegraphics[width=0.50\textwidth]{logo/WRlogo.ps}
  \label{fig:wr_logo}
\end{figure}

%\vspace{2cm}
%\begin{center}
%\large \textbf{Acknowledgements:} 
%\end{center}

%\vspace{1.3cm}

%\begin{center}
%Many thanks to John Eidson and the White Rabbit Community\\ for the feedback and comments.
%\end{center}

\newpage

\setcounter{page}{1}
\begin{center}
\large Revision History Table 
\end{center} 


\begin{table}[ht!]
%\begin{table}[tbp]
%\caption{White Rabbit Specification Revision History Table } 
\centering
\begin{tabular}{| c |  c | c | p{8cm}    |}          \hline
\textbf{Version} & \textbf{Date} & \textbf{Authors} & \textbf{Description}  \\ 
 & &  &\\ \hline
0.1 & 10/09/2010 & E.G., M.L.   & First draft for comments.  \\ \hline
0.2 & 7/09/2010  & T.W., M.L.   & Added introduction about PTP.  \\ \hline
0.3 & 8/09/2010  & J.S., M.L.   & Language\&Style-related corrections, reference to Peter's paper.  
\\ \hline
0.4 & 14/09/2010 & E.V.D.B., M.L.& 1.~Merged the FSM of the Slave and the FSM of the Master into 
                                      a single and linear FSM,\\
    &            &              & 2.~Changed WR managementIDs,\\
    &            &              & 3.~Added authors, this rev. table, explanatory figures and 
                                                              descriptions. \\ \hline
0.5 & 17/09/2010 & E.V.D.B., J.S.& 1.~Added Table of Contents. \\
    &            &               & 2.~Re-ordered appendixes. \\
    &            &               & 3.~Fixed some typos. \\ 
\hline
1.0 & 26/04/2011 & M.L.          & 1.~Included feedback from comments. \\
    &            &               & 2.~Created WR PTP Profile. \\
    &            &               & 3.~Added modification to BMC. \\ 
    &            &               & 4.~Enabled WR Switch to be a "real" Boundary Clock. \\ 
    &            &               & 5.~Added figure clarifying WRPTP vs. PTP FSMs operation from 
                                      Power On. \\ 
    &            &               & 6.~Mentioned (missing) modification to PTP FSM. \\ 
\hline
1.1 & 2/05/2011  & J.S., T.W., M.L. & Clarity and typos. \\ 
\hline
1.2 & 26/05/2011 & M.L.          & Modified according to John Eidson's feedback: . \\
    &            &               & 1.~Clear distinction between HW and WRPTP (interface defined). \\
    &            &               & 2.~Clarification of WR Message sending/receiving. \\
    &            &               & 3.~WR Link Setup in PTP Master state on the WR Master. \\
    &            &               & 4.~Added information about HW requirements and HW 
				      implementation. \\
    &            &               & 5.~WR Data Set field initialization defined and specified. \\
    &            &               & Other cosmetic changes. \\
\hline
1.3 & 7/06/2011  & M.L.          & Added Appendixes: \\
    &            &               & 1.~Extract from Tomek's MSc. \\
    &            &               & 2.~WR computation of offset and delay with asymmetry correction.\\
\hline
\end{tabular}
\label{tab:revHist}
\end{table}

\newpage

\begin{table}[ht!]
%\begin{table}[tbp]
%\caption{White Rabbit Specification Revision History Table } 
\centering
\begin{tabular}{| c |  c | c | p{8cm}    |}          \hline
\textbf{Version} & \textbf{Date} & \textbf{Authors} & \textbf{Description}  \\ 
 & & \textcolor{white}{E.V.D.B., J.S} &\\ \hline
1.4 & 02/07/2011 & M.L.          & Changes based on ptpx daemon implementation: \\
    &            &               & 1.~Name changed: REQ\_CALIBRATION to CALIBRATION.\\
    &            &               & 2.~Included Tx calibration into the CALIBRATION.\\
    &            &               & 3.~Added Link Down event description.\\
    &            &               & 4.~Specified additional wrModeOn's re-setting.\\
    &            &               & 5.~Re-modified BMC, completed DS Update tables.\\
    &            &               & 6.~Added fields to DS: calRetry and otherPortCalRetry.\\
    &            &               & 7.~Added calRetry field to WR TLV CALIBRATE Signaling Message.\\
\hline
2.0 & 04/07/2011 & J.S.          & Clarity and typos (version 2.0 release). \\
\hline
\end{tabular}
\label{tab:revHist}
\end{table}

\newpage

\tableofcontents

\newpage

\section{Introduction}

White Rabbit (WR) is a protocol developed to synchronize nodes in 
a packet-based network with sub-ns accuracy. The protocol results 
from the combination of IEEE1588-2008 (PTP)\cite{IEEE1588} with two further
requirements: precise knowledge of the link delay and clock
syntonization\footnote{The adjustment of two electronic circuits or devices in terms of frequency.} 
over the physical layer with Synchronous Ethernet (SyncE) \cite{SynchE}.

A WR link is formed by a pair of nodes, master and slave 
%(Figure~\ref{fig:wrNodes})
. The master
node uses a traceable clock to encode data over the physical layer, while the slave
recovers this clock (synchronization) and bases its timekeeping on it. Absolute time
synchronization between master and slave is achieved by adjusting
the clock phase and offset of the slave to that of the master. 
The \textit{offset} refers to the clock (e.g. time defined in Coordinated Universal Time or
International Atomic Time standards), 
while the \textit{phase} refers to the clock signal (e.g. 125 MHz clock signal). 
The phase and offset adjustment is done through the two-way exchange of PTP sync messages, which 
are time-stamped
to achieve sub-ns accuracy thanks to hardware support such as described in section~\ref{sec:hw}.
Details of a sample hardware implementation are presented in 
Appendix~\ref{sec:sampleHW}.

In White Rabbit, the precise knowledge of the link delay is obtained by accurate hardware 
timestamps (see section~\ref{sec:timestamps}) and calculation of the delay asymmetry (see 
section~\ref{sec:delayAsymCal}) supported by knowledge of delays introduced by the hardware (see 
section~\ref{sec:fixedDelays}).

The described single-link synchronization can be replicated. 
Multi-link WR networks are obtained by chaining WR links forming a
hierarchical topology. This hierarchy is imposed by the fact that
a frequency traceable to a common grandmaster must be distributed
over the physical layer, resulting in a \textit{cascade} of master
and slave nodes. As a result of this topology, a WR~network consists of two kinds of WR network 
devices:
\textit{WR boundary clocks} (WR Switches) and \textit{WR ordinary clocks} 
(WR Nodes), see section~\ref{sec:definitions}. A WR boundary clock distributes the frequency 
and time retrieved from the upstream link to all the downstream links.

\begin{figure}[ht!]
  \centering
%  \vspace{-1.3cm}
  \includegraphics[width=0.60\textwidth]{network/wrTopology.eps}
  \caption{White Rabbit network; it forms a hierarchical topology.}
  \label{fig:wrNetwork}
\end{figure}

Some applications need WR and IEEE1588-2008 nodes to coexist. Examples
of this are networks where the need for highly accurate time synchronization is concentrated on a
certain group of nodes. For this purpose the WR protocol enables WR Nodes or Switches 
to defer to IEEE1588 behavior when not connected to another WR Node or Switch.




\newpage

\section{Precision Time Protocol}
\label{sec:ptp}

The IEEE1588-2008 standard, known as Precision Time Protocol (PTP), 
is repeatedly referenced in this document. Knowledge of basic PTP concepts is 
required to read this specification. Therefore, they are explained in this section. 

PTP is a packet-based protocol designed to synchronize devices in distributed systems. 
The standard defines two kinds of messages which are exchanged between \textit{PTP nodes}: 
\textit{event messages} and \textit{general messages}. Both, the time of transmission and 
the time of reception of event messages are \textit{recorded}. General messages 
are used by PTP nodes to identify other PTP nodes, establish clock hierarchy and 
exchange data, e.g.~timestamps, settings or parameters. PTP defines several methods 
for node synchronization. Figure~\ref{fig:wrPTPmsgs} presents the messages used when 
the \textit{delay request-response mechanism} (with a \textit{two-step clock}) is used, 
which is the case in White Rabbit. 
% For simplicity reasons, a PTP node is 
% considered an \textit{ordinary clock} in the remainder of this section; such clocks 
% have only one port. 

\begin{figure}[ht!]
  \centering
  \includegraphics[width=0.30\textwidth]{protocol/ptpMSGs.ps}
  \caption{PTP messages used by WRPTP.}
  \label{fig:wrPTPmsgs}
\end{figure}

An \textit{Announce Message} is periodically broadcast by the PTP node which is 
in the Master state. The message carries information about its originator and 
the originator's clock source quality. This enables other PTP nodes receiving 
the announce message to perform the Best Master Clock (BMC) Algorithm. 
This algorithm defines the role of each PTP node in the PTP network hierarchy; 
the outcome of the algorithm is the recommended next state of the PTP node and 
the node's synchronization source (grandmaster). In other words, a PTP node 
decides to which other PTP node it should synchronize based on the information 
provided in the announce messages and using the BMC algorithm. A PTP node which is in 
the SLAVE state synchronizes to the clock of another PTP node. A PTP node 
which is in the MASTER state is regarded as a source of synchronization 
for the other PTP nodes. The full PTP state machine with state descriptions 
is included in Appendix~\ref{ptpFSM}.



\textit{Sync Messages} and \textit{Delay\_Req Messages} are timestamped 
($t_{1}$, $t_{2}$, $t_{3}$, $t_{4}$) and these timestamps are used to calculate 
the offset and the link delay between the nodes exchanging the messages. \textit{Follow\_Up Messages} 
and \textit{Delay\_Resp Messages} are used to send timestamps between Master 
and Slave (in the case of a two-step clock). 

\textit{Management Messages} are used only for configuration and 
administrative purposes. \textit{Signaling Messages} are used for communication between clocks 
for optional or experimental features of the PTP standard as well as implementation-specific 
mechanisms. Both, Management and Signaling messages, are not essential for PTP synchronization.


The flow of events in the PTP delay request-response (two-step clock) mechanism is the following 
(simplified overview): 
\begin{enumerate}
\item The master sends Announce messages periodically.
\item The slave receives the Announce message and uses the BMC algorithm to establish its place in 
      the network hierarchy.
\item The master periodically sends a Sync message (timestamped on transmission, $t_1$) followed by 
      a Follow\_Up message which carries $t_1$.
\item The slave receives the Sync message sent by the master (timestamped on reception, $t_2$).
\item The slave receives the Follow\_Up message (which carries the Sync transmission time,  
      $t_1$) sent by the master .
\item The slave sends a Delay\_Req message (timestamped on transmission, $t_3$).
\item The master receives the Delay\_Req message sent by the slave (timestamped on reception, 
      $t_4$).
\item The master sends the Delay\_Resp message which carries $t_4$.
\item The slave receives the Delay\_Resp.
\item The slave adjusts its clock using the clock offset and the link delay calculated with timestamps 
      ($t_{1}$, $t_{2}$, $t_{3}$, $t_{4}$). This results in the Slave's synchronization with the
      Master clock.
\item Repeat 1-10.
\end{enumerate}


\newpage

\section{Link Delay Model}
\label{sec:delaymodel}

The delay of a message traveling from master to slave (see Figure~\ref{fig:delaymodel}) 
can be expressed as the sum

\begin{equation}
  \label{eq:delayms}
  \eqdelay{ms} = \Delta_{tx_m} + \delta_{ms} + \Delta_{rx_s}
\end{equation}

where $\Delta_{tx_m}$ is the fixed delay due to the master's transmission
circuitry, $\delta_{ms}$ is the variable delay incurred in the
transmission medium and $\Delta_{rx_s}$ is the fixed delay due to the
slave's reception circuitry.
In a similar fashion, the delay of a message traveling from slave to
master can be decomposed as

\begin{equation}
  \label{eq:delaysm}
  \eqdelay{sm} = \Delta_{tx_s} + \delta_{sm} + \Delta_{rx_m}
\end{equation}

The characterization of the link is completed with an equation to
relate the two variable delays $\delta_{ms}$ and $\delta_{sm}$.
%  From now
% on in this document we refer to this missing equation as the \textit{physical
% medium correlation}. 
Describing a procedure to obtain this equation is out of the scope of this document. However, 
section~\ref{sec:physcorr} provides such an equation obtained empirically for one scenario.

\begin{figure}[ht!]
  \centering
    \includegraphics[width=0.95\textwidth]{protocol/delaymodel.eps}
  \caption{Delay model of a WR link. The timestamps are accurately corrected
    for link asymmetries by the usage of the four fixed delays
    $\Delta_{\{tx_m, rx_s, tx_s, rx_m\}}$ and the relationship
    between both variable delays $\delta_{\{ms, sm\}}$. }
  \label{fig:delaymodel}
\end{figure}

\subsection{Relative Delay Coefficient}
\label{sec:physcorr}

An accurate relation between both variable delays on the transmission line
is essential for obtaining an acceptable estimate of the delay asymmetry on a
WR link. The relation between $\delta_{ms}$ and $\delta_{sm}$ is represented in 
this document by the \textit{relative delay coefficient} ($\alpha$).
Its origin is highly implementation-dependent. Thus this document just assumes that it exists 
and is known. 

\subsubsection{1000base-X over a Single-mode Optical Fiber}
\label{sec:singlefiber}
When a single-mode fiber is used as bi-directional communication medium,
it can be shown that both variable delays are related by an equation of
the form \cite{Peek2010}:
\begin{equation}
\delta_{ms} = (1 + \alpha) \, \delta_{sm}
\label{eq:singlefiber}
\end{equation}

\section{Delay Asymmetry Calculation}
\label{sec:delayAsymCal}

Let us start from the PTP sync timestamps, represented by the familiar
set $t_1$, $t_2$, $t_3$ and $t_4$. The mean path delay is then defined as
\begin{equation}
\label{eq:meanPath}
  \mu = \frac{(t_2 - t_1) + (t_4 - t_3)}{2}
\end{equation}
Note that the transmission delays $t_2 - t_1$ and $t_4 - t_3$ can be
expressed in terms of WR's Delay Model:
\begin{align}
\
  t_2 - t_1 = \Delta_{tx_m} + \delta_{ms} + \Delta_{rx_s} + \eqoffset{ms}\\
  t_4 - t_3 = \Delta_{tx_s} + \delta_{sm} + \Delta_{rx_m} - \eqoffset{ms}
\end{align} where \eqoffset{ms} is the time offset between the
slave's clock and the master's. Combining the three equations above
we obtain
\begin{equation}
\label{eq:mu}
  2\mu = \Delta + \delta_{sm} + \delta_{ms}
\end{equation} where $\Delta$ accounts for all fixed delays in the path,
i.e.
\begin{equation}
\label{eq:delta}
  \Delta = \Delta_{tx_m} + \Delta_{rx_s} + \Delta_{tx_s} + \Delta_{rx_m}
\end{equation}

The delay asymmetry as specified in section 7.4.2 of the PTP standard is expressed in our own 
notation by using equations \eqref{eq:delayms},
\eqref{eq:delaysm} and \eqref{eq:mu} as follows:
\begin{align}
  \eqdelay{ms} = \mu + \eqasymm
  \label{eq:asymm} \\
  \eqdelay{sm} = \mu - \eqasymm
\end{align}


\subsection{Solution for Ethernet over a Single-mode Optical Fiber}
\label{sec:EthernetOverSingleModeFiber}
Combining equations \eqref{eq:singlefiber} and \eqref{eq:mu} we obtain:
\begin{gather}
  \delta_{ms} = \frac{1 + \alpha}{2 + \alpha} \, (2\mu - \Delta)
  \label{eq:deltams} \\
  \delta_{sm} = \frac{2\mu - \Delta}{2 + \alpha}
  \label{eq:deltasm}
\end{gather}
\\
The delay asymmetry can then be derived from equations \eqref{eq:delayms}, \eqref{eq:asymm}, 
\eqref{eq:deltams} and
\eqref{eq:deltasm}:
\begin{equation}
\label{eq:aqasymm}
  \eqasymm = \Delta_{tx_m} + \Delta_{rx_s} - \frac{\Delta - \alpha \mu + \alpha \Delta}{2 + \alpha}
\end{equation}
\\


\newpage

\section{Hardware Support}
\label{sec:hw}


This section gives a general overview of hardware requirements to support the White Rabbit protocol. 
Hardware supporting WR must fulfill the following requirements:
\begin{itemize}
  \item It shall distribute frequency over physical layer with SyncE.
  \item It shall feature constant rx/tx latencies during operation and inform higher layers 
	about these latencies.
  \item It shall provide timestamps with a sufficient precision. 
  \item It shall be able to generate the WR calibration pattern 
	(RD+K28.7 code group, Appendix 36A.2 of \cite{IEEE802.3}).
\end{itemize}

Figure~\ref{fig:clocksHwSpec} (a) gives a general picture of the data flow in the WR 
protocol, WR Hardware Support and the interface between theme
Figure~\ref{fig:clocksHwSpec} (b) explains the WR synchronization and synchronization scheme.
An example hardware implementation is described in detail in Appendix~\ref{sec:sampleHW}.

\begin{figure}[ht!]
  \centering
  \includegraphics[width=0.9\textwidth]{protocol/clocksHwSpec.ps}
  \caption{WR protocol and WR Hardware overview.}
  \label{fig:clocksHwSpec}
\end{figure}


\subsection{Synchronous Ethernet}
\label{sec:syncE}
SyncE is responsible for clock syntonization (frequency transfer) in the White Rabbit Network. 
In the SyncE scheme, the reference clock (125MHz) is used to encode the outgoing data
stream. The same clock is retrieved on the other side of the physical link.
The retrieved frequency can be further distributed and is always looped back to the sender.
Having the same frequency allows the use of phase detector technologies as a means of evaluating
delays. The measurement of phase can be used for increasing the precision of timestamps 
beyond the resolution allowed by the 125MHz clock (section~\ref{sec:timestamps}) and to obtain 
the transmit/receive (Rx/Tx) fixed delays introduced by the PHY (section~\ref{sec:fixedDelays}).

\subsection{Fixed delays}
\label{sec:fixedDelays}
The knowledge of all fixed delays $\Delta_{\{tx_m, rx_s, tx_s, rx_m\}}$ is necessary to calculate 
the delay asymmetry using the WR Link Model (section~\ref{sec:delaymodel}). They might include circuit, 
components (e.g. PHY, SFP) and FPGA internal latencies (Appendix~\ref{s:asymmetry} describes 
fixed delays in the WR optical link model). Such delays may be constant for 
the lifetime of the hardware, its up-time or the duration of the link connection. Therefore, 
the method for obtaining fixed delays is medium-specific and implementation-dependent. 
The WR protocol requires the fixed delay values and enables PHY delay
measurement, but does not specify how they shall be obtained.

The PHY's delays are measured (if necessary) and their values are distributed across the link during 
the process of establishing the WR link, which is called \textit{WR Link Setup} in this document. 
A WR~node participates in the measurement of another 
WR node's reception fixed delays ($\Delta_{\{rx_m, rx_s\}}$) upon request, e.g. by sending 
a calibration pattern in Gigabit Ethernet. The calibration pattern sent by WR nodes to measure 
fixed delays shall be generated by a repeated transmission of the RD+ K28.7 code-group. 
Measurement of fixed delays during WR Link Setup is optional and in principle needed only once, 
while the WR link is being set up. It is only required if non-deterministic reception/transmission 
elements (i.e. with latencies varying after each power cycle or lock) are used.

Figure~\ref{fig:clocksHwSpec} (a) presents a possible method of fixed delay measurement for a 
non-deterministic Gigabit Ethernet PHY -- detection of the clock signal's phase difference between 
the input and output of the PHY (detailed description in Appendix~\ref{sec:calibForGigbitE}). 


\subsection{Timestamps}
\label{sec:timestamps}

WR requires precise timestamps. 
The precision of timestamps shall be sufficient to take advantage of the calculated asymmetry
and is dictated by the required synchronization accuracy. 

In standard implementations of timestamping units the precision is limited by 
the timestamping resolution (e.g. 8ns for 125MHz). 
In particular, the precision of the reception timestamps ($t_2$, $t_4$) needs to be enhanced. 
This can be done by using the common notion of frequency distributed with SyncE to cast the 
problem of timestamping into a phase detection measurement.

% For example, in order to achieve sub-nanosecond synchronization accuracy, the precision of 
% timestamps provided by the hardware shall be in a picoseconds range. 
An overview of a possible solution enabling highly precise timestamping is illustrated in 
Figure~\ref{fig:clocksHwSpec}. The timestamping unit produces timestamps on both the rising 
and falling edges of the clock. It is provided with the phase measurement of the round-trip phase 
($phase_{MM}$) in case of the Master, and the setpoint phase ($phase_{S}$) in case of the Slave. 
The phase measurement is used by the timestamping unit to choose the correct edge timestamp and 
enhance it to the precision allowed by the phase detector. 
Details of the implementation of the presented solution are included in Appendix~\ref{s:fine_delay}.




\newpage

\section{White Rabbit PTP }


White Rabbit extends the IEEE1588-2008 (PTP) standard to achieve sub-ns accuracy while still 
benefiting from PTP's synchronization, management and messaging mechanisms. From now on in this
document, the White Rabbit extension to the PTP standard will be referred to as \textit{WRPTP}.

WRPTP takes advantage of PTP's customization facilities, i.e. \textit{PTP profiles} and 
\textit{Type-Length-Value} (TLV).
%, and makes an extension to the PTP standard which is out of the scope of these facilities. 
It also defines \textit{implementation specific} functionalities (e.g. WR-specific fields of 
Data Sets and hardware support requirements) which are compatible with the PTP standard. 
For the sake of simplicity and clarity, this section explains all the mechanisms of WRPTP without 
classifying them into PTP-profile
% , PTP-extension 
and implementation-specific. The distinction is 
made in the last subsection, where the \textit{White Rabbit PTP profile} is defined and 
the implementation-specific recommendations are listed.
%and the extensions to PTP are listed.

\subsection{Overview}
\label{sec:wrptpOverview}


WRPTP introduces \textit{the White Rabbit Link Setup} (WR Link Setup), which is a process for 
establishing the WR link  (Figure~\ref{fig:wrptpMSGs}). It includes syntonization of the local clock
 over the physical layer, measurement of fixed delays (calibration)  and distribution of the 
information about fixed delays over the link. WRPTP uses the Link Delay 
Model to obtain an accurate delay estimation, e.g. it uses the delay asymmetry equation 
\eqref{eq:aqasymm}  for Gigabit Ethernet over optical fiber. The additional communication between 
a master and a slave needed to exchange parameters and state information is done through extended 
PTP messaging facilities. 

WRPTP is compliant with standard PTP. Figure~\ref{fig:hybrid-network} depicts the tree topology 
of a \emph{hybrid} WR/IEEE1588 network with a grandmaster as a root. 

WRPTP protocol cannot be performed over a non-WR network device (i.e. switches, repeaters, routers). 
For example, if a non-WR device is connected between two WR-compatible nodes, 
the standard PTP protocol will be used for synchronization. 


\begin{figure}[ht!]
  \centering
%  \vspace{-1.3cm}
  \includegraphics[width=0.8\textwidth]{network/hybrid.ps}
  \caption{Hybrid WR/IEEE1588 network. White Rabbit nodes work
    transparently with PTP nodes. WR ordinary clock 3 is more accurately
      synchronized to the grandmaster than WR ordinary clock 2, which
      is below a standard PTP boundary clock.}
  \label{fig:hybrid-network}
\end{figure}


\begin{figure}[ht!]
  \centering
%  \vspace{-1.3cm}
  \includegraphics[width=0.99\textwidth]{protocol/wrptpMSGs_1.ps}
  \caption{Simplified overview of the message flow in WRPTP.}
  \label{fig:wrptpMSGs}
\end{figure}

\newpage 

The flow of events for standard PTP which is presented in section~\ref{sec:ptp} is extended 
as depicted in Figure~\ref{fig:wrptpMSGs} and described below:
\newpage
\begin{enumerate}
\item \textbf{WR Node A} which is in PTP\_MASTER state periodically sends WR Announce messages with 
      a custom suffix.
\item \textbf{WR Node B} receives Announce message(s), recognizes the WR Announce message and 
      uses the modified BMC algorithm (section~\ref{sec:wrBMC}) to establish its place in the WR 
      network hierarchy.
\item \textbf{WR Node B} enters WR Slave mode (based on the conditions in 
      \ref{sec:wrSlaveFSMstart}) and starts the WR Link Setup by sending the SLAVE\_PRESENT 
      WR Signaling message.
\item \textbf{WR Node A} enters WR Master mode (based on the conditions in 
      \ref{sec:wrMasterFSMstart}) and sends the LOCK WR Signaling message to request the WR Slave 
      to start syntonization.
\item The WR Slave sends the LOCKED WR Signaling message as soon as the syntonization process is 
finished (notification from the hardware).
\item The WR Master sends the CALIBRATE WR Signaling message to request the calibration pattern. 
It calibrates its transmission and reception fixed delay.
\item The WR Master sends the CALIBRATED WR Signaling message as soon as the calibration is 
finished (notification from hardware).
\item The WR Slave sends the CALIBRATE WR Signaling message to request the calibration pattern. 
It calibrates its transmission and reception fixed delay.
\item The WR Slave sends the CALIBRATED WR Signaling message as soon as the calibration is 
finished (notification from hardware).
\item The WR Master sends the WR\_MODE\_ON WR Signaling message to indicate completion of  the WR
Link Setup process.
\item The WR Master periodically sends a Sync message (timestamped on transmission, $t_1$)  followed
by a Follow\_Up message which carries $t_1$.
\item The WR Slave receives the Sync message sent by the master (timestamped on reception, $t_2$).
\item The WR Slave receives the Follow\_Up message sent by the master.
\item The WR Slave sends a Delay\_Req message (timestamped on transmission, $t_3$).
\item The WR Master receives the Delay\_Req message sent by the slave (timestamped on reception, $t_4$).
\item The WR Master sends a Delay\_Resp message which carries $t_4$.
\item The WR Slave receives the Delay\_Resp.
\item The WR Slave adjusts its clock using the clock offset and the link delay calculated with 
the timestamps ($t_{1}$, $t_{2}$, $t_{3}$, $t_{4}$). This results in the slave's synchronization 
with the master clock with sub-ns accuracy.
\item Repeat 1, 11-18.
\end{enumerate}

\newpage

\subsection{Definitions}
\label{sec:definitions}


The following definitions used in this document are based on Clause 3.1 of PTP:\\
\textbf{node}:  A device that can issue or receive Precision Time Protocol (PTP) communications on a
network.\\
\textbf{boundary clock}: A clock that has multiple Precision Time Protocol (PTP) ports in a domain 
and maintains the timescale used in the domain. It may serve as the source of time and frequency, 
i.e., be a master clock, and may synchronize and syntonize to another clock, 
i.e., be a slave clock. The frequency retrieved by the slave clock is re-distributed 
to syntonize other slave clocks connected to its ports.\\
\textbf{ordinary clock}: A clock that has a single Precision Time Protocol (PTP) port in a domain 
and maintains the timescale used in the domain. It may serve as a source of time and frequency, 
i.e., be a master clock, or may synchronize and syntonize to another clock, 
i.e., be a slave clock.\\
\\
This document introduces the following definitions:\\
\textbf{White Rabbit Port (WR Port)}: a WRPTP-enabled port of a boundary clock 
or an ordinary clock.\\
\textbf{White Rabbit Master (WR Master)}: a WRPTP-enabled port of a boundary clock or an 
ordinary clock which acts as a source of time and frequency for another White Rabbit enabled port of 
a boundary or an ordinary clock.\\
\textbf{White Rabbit Slave (WR Slave)}: a WRPTP-enabled port of a boundary clock or an 
ordinary clock which retrieves time and frequency sent over a link by 
the WR Master.\\
\textbf{White Rabbit Switch (WR Switch)}: a boundary clock implementing WRPTP, compatible with 
standard PTP. \\
\textbf{White Rabbit Node (WR Node)}: an ordinary clock implementing WRPTP, compatible with 
standard PTP. \\





\subsection{WRPTP Data Set Fields}
\label{sec:wrDS}

The PTP standard defines data sets (DS) to store the static and dynamic variables needed for  the
operation of the protocol (clause 8, PTP). WRPTP:
\begin{itemize}
  \item adds fields to the DSs defined in the PTP standard to store the WR-specific parameters,
  \item defines a new DS: backupParentDS. 
\end{itemize}


\subsubsection{New fields of PTP Data Sets}
Table~\ref{tab:wrDS} defines the additional DS fields required by WRPTP that are not part of
the PTP standard (see Appendix~\ref{ap:wrDataTypes} for a definition of primitive data types). 

\newpage

\begin{table}[ph!]
\caption{WRPTP Data Sets fields}
\centering
\begin{tabular}{| l     |     c        |           l             |   c   |   c   | c |  }        \hline   
   
\textbf{DS member}     	& \textbf{DS name}& \textbf{Allowed}       & \textbf{Initialization}& 
\textbf{D}ynamic & \textbf{Unit}            \\  
			&              & \textbf{values}           &\textbf{values}& 
or \textbf{S}tatic    &                    
                                                                                        \\ \hline
wrConfig   		& portDS       & NON\_WR,		   &       -       & S & - \\
			&	       & WR\_S\_ONLY,  		   &		   &   &   \\
			&	       & WR\_M\_ONLY, 		   &		   &   &   \\
			&	       & WR\_M\_AND\_S	           &               &   &   \\ \hline
wrMode			& portDS       & NON\_WR,\footnotemark[2]  & NON\_WR	   & D & - \\
			&	       & WR\_SLAVE,  		   &		   &   &   \\
			&	       & WR\_MASTER                &    	   &   &   \\ \hline
wrModeOn 		& portDS       & TRUE, FALSE               & FLASE         & D & - \\ \hline
wrPortState 		& portDS       & Table~\ref{tab:wrFSMdesc} & IDLE          & D & - \\ \hline
knownDeltaTx   		& portDS       & UInteger64                &       -       & S & [$2^{16}$ps]\footnotemark[3]\\ \hline
knownDeltaRx   		& portDS       & UInteger64                &       -       & S & [$2^{16}$ps]\footnotemark[3]\\ \hline
deltasKnown		& portDS       & TRUE, FALSE               &       -       & S & - \\ \hline
calibrated 		& portDS       & TRUE, FALSE               & deltasKnown   & D & - \\ \hline
deltaTx       		& portDS       & UInteger64                & knownDeltaTx  & D & [$2^{16}$ps]\footnotemark[3]\\ \hline
deltaRx     		& portDS       & UInteger64                & knownDeltaRx  & D & [$2^{16}$ps]\footnotemark[3]\\ \hline

wrStateTimeout		& portDS       & UInteger32                &       -       & S & [ms]\\ \hline
wrStateRetry		& portDS       & UInteger8                 &       -       & S & - \\ \hline

calPeriod   		& portDS       & UInteger32\footnotemark[4]&    -       & S & [us]\\ \hline
calRetry   		& portDS       & UInteger8\footnotemark[5]&    -       & S & -\\ \hline

%%%%%%%%%%%%%%%%%%%%%%%%%%%%%%%%%%%%%%%%%%%%%%%%%%%%%%%%%%%%%%%%%%%%%%%%%%%%%%%%%%%%%%%%%%
parentWrConfig 		& portDS       & NON\_WR,  		  & NON\_WR	   & D & - \\
			&	       & WR\_S\_ONLY,   	  &		   &   &   \\ 
			&	       & WR\_M\_ONLY,  		  &		   &   &   \\
			&	       & WR\_M\_AND\_S            &                &   &   \\ \hline
parentWrMode		& portDS       & NON\_WR,		  & NON\_WR 	   & D & - \\
			&	       & WR\_SLAVE,    		  &		   &   &   \\
			&	       & WR\_MASTER               &	           &   &   \\ \hline
parentWrModeOn 		& portDS       & TRUE, FALSE              & FALSE          & D & - \\ \hline
parentCalibrated	& portDS       & TRUE, FALSE              & FALSE          & D & - \\ \hline
%%%%%%%%%%%%%%%%%%%%%%%%%%%%%%%%%%%%%%%%%%%%%%%%%%%%%%%%%%%%%%%%%%%%%%%%%%%%%%%%%%%%%%%%%%
otherPortDeltaTx      	& portDS       & UInteger64               & 0x0            & D & [$2^{16}$ps]\footnotemark[2]\\ \hline
otherPortDeltaRx     	& portDS       & UInteger64               & 0x0            & D & [$2^{16}$ps]\footnotemark[2]\\ \hline
otherPortCalPeriod   	& portDS       & UInteger32               & 0x0            & D & [us]\\ \hline
otherPortCalRetry   	& portDS       & UInteger8		  & 0x0       	   & D & - \\ \hline
otherPortCalSendPattern & portDS       & TRUE, FALSE              & FALSE          & D & - \\ \hline

%%%%%%%%%%%%%%%%%%%%%%%%%%%%%%%%%%%%%%%%%%%%%%%%%%%%%%%%%%%%%%%%%%%%%%%%%%%%%%%%%%%%%%%%%%
primarySlavePortNumber	& currentDS    & UInteger16               & 0x0           & D  & - \\ \hline

\end{tabular}
\label{tab:wrDS}
\end{table}
\footnotetext[2]{Not re-initialized on exiting IDLE state of WR State Machine.}
\footnotetext[3]{The value of $\Delta_{tx_m,rx_s,rx_m,tx_s}$ measured in picoseconds and 
multiplied by $2^{16}$.}
\footnotetext[4]{Maximum shall be smaller than AnnounceInterval (see PTP).}
\footnotetext[5]{Maximum shall be 32.}
\addtocounter{footnote}{4}

\newpage

\paragraph{(Re-)Initialization}
\label{sec:initialization}

Initialization rules listed in PTP (clause 8.1.3) apply to WR fields. 
In particular:
\begin{itemize}
  \item Data set members shall be initialized before leaving the PTP INITIALIZATION state,
  \item Static members shall be initialized to the implementation-specific values meeting the 
        specifications for the member.
\end{itemize}
Additionally, the values of dynamic WR fields shall be set to initializing values on the occurrence 
of the following events:
\begin{itemize}
  \item exiting IDLE state of the WR State Machine (except for wrMode).
  \item EXCEED TIMEOUT RETRIES transition event.
\end{itemize}
The initialization values of dynamic WR fields are specified in Table~\ref{tab:wrDS}, while 
the initialization values of static fields are implementation and configuration-specific 
(if not defined by implementation, the default values in Table~\ref{tab:wrDSdefault} shall be used).



\begin{table}[ph!]
\caption{Default values for static DS.}
\centering
\begin{tabular}{| l     |     c           |           c           |    }        \hline   
   
\textbf{DS member}     	& \textbf{DS name}& \textbf{Default value} \\   
			&                 & 			   \\ \hline
wrConfig   		& portDS          & WR\_M\_AND\_S	   \\ \hline
deltasKnown		& portDS          & FALSE                 \\ \hline
knownDeltaTx   		& portDS          & 0 [$2^{16}$ps]\footnotemark[2]                 \\ \hline
knownDeltaRx   		& portDS          & 0 [$2^{16}$ps]\footnotemark[2]                 \\ \hline
calPeriod   		& portDS          & 3000 [us]              \\ \hline
calRetry   		& portDS          & port number + 2         \\ \hline
wrStateTimeout		& portDS          & 1000 [ms]              \\ \hline
wrStateRetry		& portDS          & 3                      \\ \hline

\end{tabular}
\label{tab:wrDSdefault}
\end{table}

\paragraph{New Fields Description}

\subparagraph{wrConfig} Determines the predefined function of a WR Port: 
\begin{itemize}
  \item WR\_M\_AND\_S shall be the default on boundary and non-slave-only ordinary clocks.
  \item WR\_SLAVE shall be the default on a WR Node which is a slave-only ordinary clock.
  \item WR\_MASTER shall be default on a WR Node (ordinary clock) which is, by design, 
	supposed to be always connected to a primary source of time, so-called master-only.
  \item NON\_WR might be used on any kind of clock to disable the WRPTP protocol on a given port. 
	Such configuration can result from administrative constraints or requirements.  
\end{itemize}

\subparagraph{wrMode} 
\label{par:wrMode}
Contains the current WR-role of a WR Port which is determined using the modified 
BMC algorithm (section~\ref{sec:wrBMC}) by the conditions described in 
sections~\ref{sec:wrSlaveFSMstart}~and~\ref{sec:wrMasterFSMstart}. 
The WR-role is recommended or active depending on the value of wrModeOn:
\begin{itemize}
  \item wrMode is \textit{recommended} when wrModeOn is FALSE.
  \item wrMode is \textit{active} when wrModeOn is TRUE.
\end{itemize}
A WR Port in WR\_SLAVE mode is called a WR Slave. Similarly, a WR Port in WR\_MASTER mode is 
called WR Master. 
The wrMode field shall be reset to NON\_WR on the reception of HW\_LINK\_DOWN hardware 
event notification (see \ref{sec:hwOutput}) and on reception of an 
ANNOUNCE\_RECEIPT\_TIMEOUT\_EXPIRES event on the WR Slave.

\subparagraph{wrModeOn} 
\label{par:wrModeON}
If TRUE, it indicates that the WR Link Setup has been performed successfully 
and the WR Port is synchronized using WRPTP, consequently the role defined in wrMode is active.

It is set to TRUE on successful completion of the WR State Machine (WR\_LINK\_ON state, 
section~\ref{wrFSM}).

It is set to FALSE :
\begin{itemize}
  \item when link-down is detected on the WR Port, section~\ref{sec:hwOutput},
  \item on initialization and re-initialization as described in section~\ref{sec:initialization},
  \item on occurrence of the WR\_MODE\_OFF event, section~\ref{sec:reestablishingWRLink}.
  \item on exiting the PTP\_SLAVE state.
\end{itemize}


\subparagraph{wrPortState} 
Stores the current state of the WRPTP state machine (see section~\ref{wrFSM}).

\subparagraph{knownDeltaTx} 
If the transmission fixed delay ($\Delta_{tx}$) is not known a priori (i.e. a non-deterministic 
PHY is used, see \ref{sec:fixedDelays}), it shall be set to 0x0. 
Otherwise, it stores the value of $\Delta_{tx}$ measured in picoseconds and multiplied 
by ${2^{16}}$.

\subparagraph{knownDeltaRx} 
If the reception fixed delay ($\Delta_{rx}$) is not known a priori (i.e. a non-deterministic 
PHY is used, see \ref{sec:fixedDelays}), it shall be set to 0x0. 
Otherwise, it stores the value of $\Delta_{rx}$ measured in picoseconds and multiplied 
by ${2^{16}}$.

\subparagraph{deltasKnown} 
TRUE if the fixed delays are known a priori (a deterministic PHY is used) and their values are 
provided in the knownDeltaTx and knownDeltaRx Data Set fields. Otherwise FALSE.


\subparagraph{calibrated}
\label{par:calibrated} 
If TRUE, it indicates that the fixed delays of the given port are currently known (measured and 
stored in deltaTx/deltaRx DS fields).

\subparagraph{deltaTx} 
\label{par:deltaTx}
Port's $\Delta_{tx}$ measured in picoseconds and multiplied by ${2^{16}}$. 
% If knownDeltaTx is not 
% equal 0x0, the value shall be set to knownDeltaRx during initialization 
% (section~\ref{sec:initialization}). Otherwise, it shall be measured during WR Link Setup.

\subparagraph{deltaRx} 
\label{par:deltaRx}
Port's $\Delta_{rx}$ measured in picoseconds and multiplied by ${2^{16}}$. 
% If knownDeltaRx is not 
% equal 0x0, the value shall be set to knownDeltaRx during initialization 
% (section~\ref{sec:initialization}). Otherwise, it shall be measured during WR Link Setup.

\subparagraph{wrStateTimeout}
\label{par:wrStateTimeout} 
Determines the timeout (in microseconds) for an execution of a state of the WR State Machine.

\subparagraph{wrStateRetry} 
\label{par:wrStateRetry}
Determines the default number of times a state of WR State Machine is re-entered 
(as a consequence of wrStateTimeout expiration) before the WR Link Setup is abandoned. If the 
number of the given state execution retries equals wrStateRetry, the EXC\_TIMEOUT\_RETRY event is 
generated (see \ref{sec:wrEventsAndConditions}). 

\subparagraph{calPeriod} 
\label{par:calPeriod}
Calibration period in microseconds required to perform the measurement of the deltaTx and deltaRx 
fixed delays (see \ref{sec:fixedDelays}). Its value shall be less than the 
value of the AnnounceInterval of PTP determined by the portDS.logAnnounceInterval 
(see clause 15.5.3.7.2.1 of PTP). If its value is greater than 0x0, it overwrites wrStateTimeout 
for the \textit{CALIBRATION} state of the WR State Machine. 
It is distributed to the partner-port to determine the timeout of its \textit{RESP\_CALIB\_REQ} 
state.

\subparagraph{calRetry} 
\label{par:calRetry}
If greater then 0x0, it overwrites wrStateRetry (\ref{par:wrStateRetry}) for 
the \textit{CALIBRATION} state of the WR State Machine. It is distributed to 
the partner port and used to overwrite wrStateRetry for the \textit{RESP\_CALIB\_REQ} state. 
The maximum allowed value is 32.

\subparagraph{primarySlavePortNumber} 
Zero value indicates that no Primary Slave is selected. 1, 2 ,...N -- values indicate the
portNumber (clause 7.5.2.3 PTP) selected as the Primary Slave, see section~\ref{sec:wrBMC}.

\subparagraph{parentWrConfig} 
Stores the value of wrConfig of the parent port which is sent in the Announce Message 
(section~\ref{sec:wrAnnounceMSG}).

\subparagraph{parentWrMode} 
Stores the value of wrMode of the parent port which is sent in the Announce Message 
(section~\ref{sec:wrAnnounceMSG}).

\subparagraph{parentWrModeOn} 
\label{par:parentWrModeON}
Stores the value of wrModeOn of the parent port which is sent in the Announce Message 
(section~\ref{sec:wrAnnounceMSG}). It shall be set to TRUE in the WR\_LINK\_ON state on the WR Slave.

\subparagraph{parentCalibrated} 
\label{par:parentCalibrated}
Stores the value of calibrated of the link-parent port which is sent in the Announce Message 
(section~\ref{sec:wrAnnounceMSG}).

\subparagraph{otherPortDeltaTx} 
\label{par:otherPortDeltaTx}
Stores the value of deltaTx of the link-partner port exchanged during the WR Link Setup. It is 
sent in the WRPTP Signaling Message (section~\ref{wrSignalingMSG}).

\subparagraph{otherPortDeltaRx} 
\label{par:otherPortDeltaRx}
Stores the value of deltaRx of the partner port exchanged during the WR Link Setup. It is 
sent in the WRPTP Signaling Message (section~\ref{wrSignalingMSG}).

\subparagraph{otherPortCalPeriod} 
\label{par:otherPortCalPeriod}
Stores the value of the calPeriod parameter of the link-partner port which is 
sent in the CALIBRATE message (\ref{par:CALIBRATE}). 
If greater than 0, the value defines the time (in microseconds) for which the calibration pattern 
should be sent by the port which received the message -- it overwrites wrStateTimeout 
(\ref{par:wrStateTimeout}) for the \textit{RESP\_CALIB\_REQ} state of the WR State Machine 
on the receiving port. It is exchanged 
during the WR Link Setup by the WRPTP Signaling Message (section~\ref{wrSignalingMSG}).

\subparagraph{otherPortCalRetry} 
\label{par:otherPortCalRetry}
Stores the value of the calRetry parameter of the link-partner port which is 
sent in the CALIBRATE message (\ref{par:CALIBRATE}). 
If greater than 0x0, it overwrites the wrStateRetry (\ref{par:wrStateRetry}) for 
the \textit{RESP\_CALIB\_REQ} state of the WR State Machine on the receiving port.


\subparagraph{otherPortCalSendPattern} 
\label{par:otherPortCalSendPattern}
Stores the value of calSendPattern sent by the WRPTP Signaling Message 
(section~\ref{wrSignalingMSG}) which indicates whether sending of the calibration pattern is 
required.

%\paragraph{parentCalPattern} Stores the value of calPattern parameter of the parent port.

%\paragraph{parentCalPatternLen} Stores the value of calPatternLen parameter of the parent port.


\subsubsection{backupParentDS data set specifications}

A boundary clock shall maintain an implementation-specific backupParentDS Data Set for the purpose 
of qualifying redundant sources of synchronization, i.e.:
\begin{itemize}
  \item backup paths to the current grandmaster clock, or
  \item paths to alternative grandmaster clock(s).
\end{itemize}
Each entry of the data set contains:

\begin{itemize}
  \item backupParentDS.secondarySlavePortNumber - the number of the port on which the Announce 
	message from a backup parent has been received (Secondary Slave port),
  \item backupParentDS.backupParentSourcePortIdentity - the portIdentity of the port on the backup 
        master.
\end{itemize}

The implementation-specific backupParentDS set shall have a minimum capacity of three backup parent 
records. The order of the records is established using the Data Set Comparison Algorithm as described
in section~\ref{StadeDecisionAlg}.

The initialization rules stated in PTP (clause 8.1.3) apply to the backupParentDS Data Set.

%on $E_{rbest}$ 
%which are qualified as backup Masters by modified State Decision Algorithm, record number 0 being the
%best (described in the section~\ref{StadeDecisionAlg}).



\subsection{Modified Best Master Clock Algorithm}
\label{sec:wrBMC}

\subsubsection{Overview}
The Best Master Clock (BMC) algorithm is used in PTP to compare clocks (determine which 
clock is the \textit{best}) and to recommend the next state of the PTP state machine (Clause 9.3 
of PTP). The BMC is used to prune the physical network topology to obtain logical tree(s) with the
\textit{best} clock, the grandmaster (preferably the one synchronized to a primary reference time 
source), as a root. In the case when more than one clock in a network is synchronized to a primary 
reference time source, the BMC produces disjoint logic trees. Each tree has a single grandmaster. 

As a result of BMC, no more than one port of a Boundary Clock can be in the PTP\_SLAVE state. Such a
solution allows for redundancy of the time source and topology but is not optimal for the continuity 
of the synchronization. For example, in case of a failure of one of the \textit{best} clocks
(reference-synchronized grandmasters), all the slave nodes in its logic tree need to switch 
synchronization to the alternate grandmaster. This requires the restart of the estimation of the 
clock drift, mean path delay and offset which might cause fluctuations in the notion of time. 

The modified BMC allows for more than one \textit{best} clock in a single domain, enabling the 
creation of a logic topology with multiple roots. A Boundary Clock running the modified BMC can have 
more than one port in the PTP\_SLAVE state. This means that timing information is exchanged between 
the Boundary Clock and more than one source of time (Ordinary Clock or Boundary Clock). At any time 
any of these links can be used to perform synchronization, including a weighted average from all 
SLAVE ports as mentioned in \cite{Takahide}.

The BMC algorithm includes the State Decision Algorithm (SDA) and the Data Set Comparison
Algorithm (DCA) and defines which fields of Data Sets should be updated depending on the outcome of
the SDA. The modifications to the BMC are detailed below. The DCA is not modified, therefore it is not
discussed.

\subsubsection{State Decision Algorithm (SDA)}
\label{StadeDecisionAlg}

The original SDA is depicted in Figure~26 of the PTP standard. The modified SDA, depicted in
Figure~\ref{fig:modifiedSDA}, enforces the $BMC\_SLAVE$ state instead of the $BMC\_PASSIVE$ state 
on the clocks with Class field value greater than 127 (block-13 of Figure~\ref{fig:modifiedSDA}). 
A port being in a $PTP\_SLAVE$ state as a result of the SDA modification (block-13) is considered 
a Secondary SLAVE port. The port which enters 
the $PTP\_SLAVE$ state based on the decision at block-11, is considered the Primary SLAVE port.
This means that:
\begin{itemize}
  \item the Primary SLAVE port and the Secondary SLAVE port(s) are synchronized to the same 
	grandmaster clock but the connection of the Primary SLAVE port is possibly better by path, or
  \item the Primary SLAVE port and the Secondary SLAVE port(s) are synchronized to different 
	grandmaster clocks of the same clockClass range (1-127 or 128-255,see Table~5 of PTP) 
	but the grandmaster clock connected to the Primary SLAVE port is better or better by path.
\end{itemize}

The best qualified Announce messages ($E_{rbest}$, see Clause 9.3.2.3 of PTP) from all Secondary 
SLAVE ports shall be compared with the DCA to determine the "second best master" and the lower order
masters. The sequence of events is as follows (the idea originated from \cite{Takahide}):
\begin{itemize}
  \item The best master is computed with the BMC -- the Primary SLAVE port established (block-11 of
	Figure~\ref{fig:modifiedSDA}) and the parentDS data set updated (S1). The port number is 
	stored in currentDS.primarySlavePortNumber.
  \item If the recommended state is $BMC\_SLAVE$ at block-13, the master is considered
        a second best master and its data is stored in $E_{srbest}$.
  \item The BMC algorithm is executed on all the $E_{srbest}$.
  \item The best master selected is considered second best master and its reception port is 
	considered the best Secondary SLAVE port.
  \item The BMC is repeated excluding $E_{srbest}$ of the best master until the set of $E_{srbest}$ 
	is empty.
\end{itemize}
The order of second and lower level best masters determined in this way is used to update the
backupParentDS data set.


\begin{figure}[ht!]
  \centering
  \includegraphics[width=0.85\textwidth]{protocol/SDA.ps}
  \caption{Modified State Decision Algorithm (modifications in red).}
  \label{fig:modifiedSDA}
\end{figure}

\newpage

\subsubsection{Update of Data Sets}

The modifications to the rules governing the update of the data sets are twofold:
\begin{itemize}
  \item Modifications to already existing "update tables" (Table~13 and Table~16, clause 9.3 of
	PTP) to accommodate the new DS fields.
  \item Addition of a new "update table" to accommodate the new backupParentDS data set and 
	the modifications in the SDA.
\end{itemize}

\begin{table}[tbp]
\caption{Modification to Table~13 of IEEE1588-2008: Update for state decision code M1 and M2.}
\centering
\begin{tabular}{| p{7.5cm} | p{7.5cm}|}           \hline
\textbf{Update this field}  		&  
\textbf{From the indicated field of the defaultDS data set of the clock unless otherwise stated} \\ 
					&  						\\ \hline
\multicolumn{2}{|c|}{currentDS}								\\ \hline
currentDS.primarySlavePortNumber   	& set to 0					\\ \hline
\multicolumn{2}{|c|}{portDS}	   	 						\\ \hline
portDS.parentWrConfig   		& portDS.wrConfig				\\ \hline
portDS.parentWrMode   			& portDS.wrMode					\\ \hline
portDS.parentWrModeOn   		& portDS.wrModeOn				\\ \hline
portDS.parentCalibrated		   	& portDS.calibrated				\\ \hline

\end{tabular}
\label{tab:modifiedM1M2}
\end{table}

\begin{table}[tbp]
\caption{Modification to Table~16 of IEEE1588-2008: Update for state decision code S1.}
\centering
\begin{tabular}{| p{7.5cm} | p{7.5cm}|}           \hline
\textbf{Update this field}  		&  \textbf{From the indicated source}  	\\ 
					&  						\\ \hline
\multicolumn{2}{|c|}{currentDS}    								\\ \hline
currentDS.primarySlavePortNumber   	& portDS.portIdentity.portNumber		\\ \hline

\end{tabular}
\label{tab:modifiedS1}
\end{table}

\begin{table}[tbp]
\caption{Update for state decision code S2}
\centering
\begin{tabular}{| p{8.5cm} | p{6.5cm}|}          \hline
\textbf{Update this field}  			&  \textbf{From the indicated source}  	\\ 
						&  					\\ \hline
\multicolumn{2}{|c|}{backupParentDS} 							\\ \hline
backupParentDS.secondarySlavePortNumber 	& portDS.portIdentity.portNumber	\\ \hline
backupParentDS.backupParentSourcePortIdentity 	& sourcePortIdentity of $E_{rbest}$	\\ \hline

\end{tabular}
\label{tab:modifiedS2}
\end{table}

\newpage

\subsection{WRPTP Messages}
\label{sec:wrMSG}

White Rabbit benefits from PTP's messaging facilities. It defines a \textit{WR Type-Length-Value} 
(WR TLV) extension to exchange WR-specific information and uses the two-step clock delay 
request-response mechanism for synchronization. In particular, it adds a suffix to the Announce 
message to enable recognition of WR nodes and uses Signaling Messages to exchange WR-specific 
information (Figure~\ref{fig:wrptpMSGs}). During the exchange of timestamps, their
fractional nanoseconds part is always included in the messages as defined in the PTP standard 
(cause 9.5.10) and explained in Appendix~\ref{ap:computationsExplanation}.

A WR port in the $PTP\_MASTER$ state announces its presence by adding a suffix to the Announce 
message. The suffix is defined by the PTP standard as a set of TLV entities (section 13.4, PTP). 
Unrecognized by standard PTP nodes WR TLVs are ignored(section 14.1, PTP), 
but read and interpreted by White Rabbit nodes. 
The information provided in the WRPTP Announce message is sufficient for another WR port to decide 
whether the WR link can be established and maintained. A WR port receiving a 
WRPTP Announce message enters (if conditions in \ref{sec:wrSlaveFSMstart} are fulfilled) 
the WR Slave mode and starts the Link Setup process (performed in the $PTP\_UNCALIBRATED$ state, 
see Appendix~\ref{ap:ptpAndWrFSMS}). It requests the sender of the WRPTP announce message 
to enter the WR Master mode (see section \ref{sec:wrMasterFSMstart}) and start 
the Link Setup process as well. During the WR Link Setup, communication between the WR Master and 
the WR Slave is performed using PTP Signaling Messages carrying WR TLVs 
(section~\ref{sec:communicationWRfsm}). Once the WR link has been established, the WR nodes use 
a PTP delay request-response mechanism (section 11.3, PTP).


\subsubsection{WR Type-Length-Value}
\label{sec:wrTLVtype}

All PTP messages can be extended by means of a standard \textit{type, length, value} (TLV) 
extension mechanism. White Rabbit uses the \textit{ORGANIZATION\_EXTENSION} TLV type \\
(tlvType=0x0003) which is defined in clause 14.3 of PTP. TLVs represented by this type 
are designated for vendors and standard organizations to extend the protocol for specific needs. 
The organization-specific TLV fields for the WR extension are defined and described in 
Table~\ref{tab:wrTLVtype}. 

\begin{table}[h!]
\caption{Organization specific TLV fields for WR (see Table~35, PTP).}
\centering
\begin{tabular}{| c | c | c | c | c | c | c | c | c | c |c|  p{4.3cm}  |}
\hline
\multicolumn{8}{|c|}{\textbf{Bits}} & \textbf{Octets} & \textbf{TLV} & \textbf{WR TLV} & \\
%\hline
\cline{0-7}
7 & 6 & 5 & 4 & 3 & 2 & 1 & 0 &  &  \textbf{Offset}  & \textbf{Content} & \textbf{Description} \\
\hline
\multicolumn{8}{|c|}{tlvType            } & 2 & 0  & 0x0003   & Organization extension, 
								see Table~34 PTP. \\ \hline
\multicolumn{8}{|c|}{lengthField        } & 2 & 2  & 8+N      & N is an even number of wrDataField
								octets. \\ \hline
\multicolumn{8}{|c|}{OrganizatinId      } & 3 & 4  & 0x080030 & OUI owned by CERN. \\
\hline
\multicolumn{8}{|c|}{organizationSubType} & 2 & 7  & magicNumber & 0xABCD - identifies WRPTP 
						      within protocols identified by CERN's OUI. \\
\cline{9-12} %\hline
\multicolumn{8}{|c|}{}                    & 1 & 9  & versionNumber& 0x01 - WRPTP version.\\
 \hline
\multicolumn{8}{|c|}{DataField}           & 2 & 10 & wrMessageID  & WR-specific messages identifier,
								    see Table~\ref{tab:wrMessageId}. \\
\cline{9-12} %\hline 
\multicolumn{8}{|c|}{  }                 & N & 12 & wrDataField   & Content of WR-specific message. \\
\hline
\end{tabular}
\label{tab:wrTLVtype}
\end{table}
The WR-specific TLVs are identified by wrMessageID which is defined in 
Table~\ref{tab:wrMessageId}. The different types of WR messages are described in subsequent 
sections.

\begin{table}[tbp]
\caption{White Rabbit Message ID values}
\centering
\begin{tabular}{| l | p{2.5cm} | c | c |}          \hline
\textbf{WR Message Name}  &  \textbf{wrMessageId value (hex)} &  \textbf{Sent in message type} \\ 
& &  \\ \hline
SLAVE\_PRESENT     &  0x1000 & Signaling \\ \hline
LOCK               &  0x1001 & Signaling \\ \hline
LOCKED             &  0x1002 & Signaling \\ \hline
CALIBRATE          &  0x1003 & Signaling \\ \hline
CALIBRATED         &  0x1004 & Signaling \\ \hline
WR\_MODE\_ON       &  0x1005 & Signaling \\ \hline
ANN\_SUFIX	   &  0x2000 & Announce  \\ \hline
\end{tabular}
\label{tab:wrMessageId}
\end{table}

\newpage

\subsubsection{WRPTP Announce Message}
\label{sec:wrAnnounceMSG}


The standard PTP Announce Message is suffixed by one WR TLV entity. 
The WRPTP Announce message has the structure defined in Table~\ref{tab:wrAnnounce}. The $dataField$
of the suffix WR TLV stores the $wrFlags$ defined in Table~\ref{tab:wrFlags}.

\begin{table}[ht!]
\caption{White Rabbit Announce Message.}
\centering
\begin{tabular}{| c | c | c | c | c | c | c | c | c | c | p{4.5cm} |}
\hline
\multicolumn{8}{|c|}{\textbf{Bits}} & \textbf{Octets} & \textbf{TLV}  & \textbf{Content} \\
%\hline
\cline{0-7}
7 & 6 & 5 & 4 & 3 & 2 & 1 & 0 &     & \textbf{Offset}  &  \\
\hline
\multicolumn{8}{|c|}{header        }  & 34 & 0  & section 13.3, PTP. \\ \hline
\multicolumn{8}{|c|}{body          }  & 30 & 34 & section 13.5, PTP.         \\ \hline
\multicolumn{8}{|c|}{tlvType       }  & 2  & 64 & 0x0003, see \ref{sec:wrTLVtype}. \\ \hline
\multicolumn{8}{|c|}{lengthField   }  & 2  & 66 & 0xA.   \\ \hline
\multicolumn{8}{|c|}{OrganizationId}  & 3  & 68 & 0x080030.               \\ \hline
\multicolumn{8}{|c|}{magicNumber   }  & 2  & 71 & 0xABCD.               \\ \hline
\multicolumn{8}{|c|}{versionNumber }  & 1  & 73 & 0x01.               \\ \hline
\multicolumn{8}{|c|}{wrMessageId   }  & 2  & 74 & 0x0010.               \\ \hline
\multicolumn{8}{|c|}{wrFlags       }  & 2  & 76 & see Table~\ref{tab:wrFlags}.    \\ \hline
\end{tabular}
\label{tab:wrAnnounce}
\end{table}


\begin{table}[ht!]
\caption{White Rabbit flags.}
\centering
\begin{tabular}{| c | c | c | c | p{5.3cm} |}          \hline

\textbf{Octet} &\textbf{Bit} & \textbf{Message type} & \textbf{Name}  & \textbf{Description}  \\ 
               &             &                       &                &              \\ \hline
0              &0   	& \multirow{2}{*}{Announce} 	& \multirow{2}{*}{wrConfig}  &
\multirow{2}{5.3cm}{The value of wrConfig parameter from portDS.} \\ \cline{1-2}
0              &1            &				& 		&\\ \hline
0              &2            & Announce              & calibrated     & TRUE if the source port 
									is calibrated.\\
									\hline
0              &3            & Announce              & wrModeOn     & The wrModeOn value 
								      (\ref{par:wrModeON}) of the 
								      sending port which shall be
								      stored in parentWrModeOn
								      (\ref{par:parentWrModeON})
								      of the receiving port.
% 								      The value of wrModeON
% 								      parameter of portDS.
								      \\ \hline
\end{tabular}
\label{tab:wrFlags}
\end{table}

\newpage


\subsubsection{WRPTP Signaling Messages}
\label{wrSignalingMSG}

White Rabbit uses Signaling Messages (defined in clause 13.12 of PTP) to exchange WR-specific 
information, except for the flags transported in the suffix of the Announce message. 

The Signaling Messages conform to the format presented in Table~\ref{tab:wrSignalingMSG}. Each WR
Signaling Message transports a single WR TLV structure as defined in Section~\ref{sec:wrTLVtype}.
The \textit{targetPortIdentity} shall be always set to the \textit{clockIdentity} of the port 
on the other side of the link (conveyed in the Announce Message). 

The Signaling Messages are used in WRPTP to trigger transitions in the WR State Machine and exchange 
WR-specific parameters. Each message is sent on entering particular states of the WR State Machine as 
described in section~\ref{sec:communicationWRfsm}.

The distinction between WR Signaling Messages is made by the \textit{wrMessageId} field of the 
\textit{WR TLV}, defined in Table~\ref{tab:wrMessageId}. Signaling Messages are exchanged only 
within a single link connection (no forwarding). This is ensured by setting the targetPortId 
appropriately. The rest of this subsection describes the WRPTP Signaling Messages in detail.

\begin{table}[h!]
\caption{PTP Signaling message fields (Table 33, PTP) }
\centering
\begin{tabular}{| c | c | c | c | c | c | c | c | c | c |}
\hline
\multicolumn{8}{|c|}{\textbf{Bits}}               & \textbf{Octets} & \textbf{TLV}  \\
%\hline
\cline{0-7}
7 & 6 & 5 & 4 & 3 & 2 & 1 & 0                     &                 & \textbf{Offset}    \\
\hline
\multicolumn{8}{|c|}{header                    }  & 34             & 0             \\ \hline
\multicolumn{8}{|c|}{targetPortIdentity        }  & 10             & 34            \\ \hline
\multicolumn{8}{|c|}{WR TLV                    }  & M              & 44             \\ \hline
\end{tabular}
\label{tab:wrSignalingMSG}
\end{table}


\paragraph{SLAVE\_PRESENT} Message sent by the WR Port which became WR Slave (entered WR Slave mode) 
to the link-partner WR Port. It initiates the WR Master mode (the port becomes WR Master) 
on the link-partner WR Port which starts the WR Link Setup. The message shall have the form specified 
in Table~\ref{tab:wrOtherTLV}.

\paragraph{LOCK} Message sent by the WR Master to the WR Slave to request 
the start of frequency locking. The message shall have the form specified in Table~\ref{tab:wrOtherTLV}.

\paragraph{LOCKED} Message sent by the WR Slave to the WR Master. It indicates the successful 
completion of frequency locking. The message shall have the format specified in 
Table~\ref{tab:wrOtherTLV}.


\begin{table}[h!]
\caption{WR TLV for SLAVE\_PRESENT, LOCK, LOCKED AND WR\_MODE\_ON Signaling Messages.}
\centering
\begin{tabular}{| c | c | c | c | c | c | c | c | c | c | p{4.5cm} |}
\hline
\multicolumn{8}{|c|}{\textbf{Bits}} & \textbf{Octets} & \textbf{TLV}  & \textbf{Content} \\
%\hline
\cline{0-7}
7 & 6 & 5 & 4 & 3 & 2 & 1 & 0 &     & \textbf{Offset}  &  \\
\hline
\multicolumn{8}{|c|}{tlvType       }  & 2  & 0  & 0x0003, see \ref{sec:wrTLVtype}. \\ \hline
\multicolumn{8}{|c|}{lengthField   }  & 2  & 2  & 0x8.   \\ \hline
\multicolumn{8}{|c|}{OrganizationId}  & 3  & 4  & 0x080030.               \\ \hline
\multicolumn{8}{|c|}{magicNumber   }  & 2  & 7  & 0xABCD.               \\ \hline
\multicolumn{8}{|c|}{versionNumber }  & 1  & 9  & 0x01.               \\ \hline
\multicolumn{8}{|c|}{wrMessageId   }  & 2  & 10 & Defined in Table~\ref{tab:wrMessageId}.\\ \hline
\end{tabular}
\label{tab:wrOtherTLV}
\end{table}



\paragraph{CALIBRATE}
\label{par:CALIBRATE}
Message sent by the WR port entering the \textit{CALIBRATION} state  (see
section~\ref{wrFSM}). It informs the other port 
whether sending a calibration pattern (see section~\ref{sec:fixedDelays}) is required (defined by
the value of \textit{calSendPattern} flag). If calibration is required, it carries  
the calibration period. The message format and parameters are described  in
Table~\ref{tab:wrCalibrateTLV}.


\begin{table}[h!]
\caption{WR TLV CALIBRATE Signaling Message }
\centering
\begin{tabular}{| c | c | c | c | c | c | c | c | c | c | p{7.5cm} |}
\hline
\multicolumn{8}{|c|}{\textbf{Bits}} & \textbf{Octets} & \textbf{TLV}  & \textbf{Content} \\
%\hline
\cline{0-7}
7 & 6 & 5 & 4 & 3 & 2 & 1 & 0 &     & \textbf{Offset}  &  \\
\hline
\multicolumn{8}{|c|}{tlvType               }  & 2  & 0  & 0x0003, see \ref{sec:wrTLVtype}. \\ \hline
\multicolumn{8}{|c|}{lengthField           }  & 2  & 2  & 0x0E.   \\ \hline
\multicolumn{8}{|c|}{OrganizationId        }  & 3  & 4  & 0x080030.               \\ \hline
\multicolumn{8}{|c|}{magicNumber           }  & 2  & 7  & 0xABCD.               \\ \hline
\multicolumn{8}{|c|}{versionNumber         }  & 1  & 9  & 0x01.               \\ \hline
\multicolumn{8}{|c|}{wrMessageId           }  & 2  & 10 & CALIBRATE.               \\ \hline
\multicolumn{8}{|c|}{calSendPattern}          & 1  & 12 & The value determines whether the calibration 
						          pattern should be sent. If the value
						 	  is 0x1, the calibration pattern is sent. 
							  If the value is 0x0, the calibration 
							  pattern is not sent. The value shall be 
							  based on the information provided by the
							  $calibrated$ Data Set field 
							  (\ref{par:calibrated}) It shall be
							  stored in otherPortCalSendPattern 
							  (\ref{par:otherPortCalSendPattern}).
							  \\ \hline
\multicolumn{8}{|c|}{calRetry (UInteger8)}  & 1 & 13  & The calRetry value 
							  (\ref{par:calRetry})
							  of the sending port which shall be stored 
							  in otherPortCalRetry 
							  (\ref{par:otherPortCalRetry}) of 
							  the receiving port.\\ \hline
\multicolumn{8}{|c|}{calPeriod (UInteger32)}  & 4 & 14  & The calPeriod value 
							  (\ref{par:calPeriod})
							  of the sending port which shall be stored 
							  in otherPortCalPeriod 
							  (\ref{par:otherPortCalPeriod}) of 
							  the receiving port.\\ \hline
% \multicolumn{8}{|c|}{calibrationPattern}    & 4 & 18 & The  value defines the calibration pattern
% 							which should be sent by the receiving
% 							node.\\ \hline
% \multicolumn{8}{|c|}{calibrationPatternLen} & 2 & 22 & The  value defines the number of bits of
% 							\textit{calibrationPattern} field which
% 							should be used as repeated pattern (starting
% 							  with the LSB). \\\hline
\end{tabular}
\label{tab:wrCalibrateTLV}
\end{table}


\newpage



\paragraph{CALIBRATED} Message sent by a WR port entering the \textit{CALIBRATED} state. 
If preceded by the \textit{CALIBRATE} message with \textit{calSendPattern} set to TRUE, 
it indicates successful completion of the calibration. The message provides the other port with the 
values of its fixed delays ($\Delta_{tx}$ and $\Delta_{rx}$). 
The message shall have the format specified in Table~\ref{tab:wrCalibratedTLV}.

\begin{table}[tbp]
\caption{WR TLV for CALIBRATED Signaling Message.}
\centering
\begin{tabular}{| c | c | c | c | c | c | c | c | c | c |  p{7.5cm}  |}
\hline
\multicolumn{8}{|c|}{\textbf{Bits}} & \textbf{Octets} & \textbf{TLV} &\\
%\hline
\cline{0-7}
7 & 6 & 5 & 4 & 3 & 2 & 1 & 0 &  &  \textbf{Offset}  & \textbf{Content}  \\
\hline
\multicolumn{8}{|c|}{tlvType               }  & 2  & 0  & 0x0003, see \ref{sec:wrTLVtype}. \\ \hline
\multicolumn{8}{|c|}{lengthField           }  & 2  & 2  & 0x18.   \\ \hline
\multicolumn{8}{|c|}{OrganizationId        }  & 3  & 4  & 0x080030.               \\ \hline
\multicolumn{8}{|c|}{magicNumber           }  & 2  & 7  & 0xABCD.               \\ \hline
\multicolumn{8}{|c|}{versionNumber         }  & 1  & 9  & 0x01.               \\ \hline
\multicolumn{8}{|c|}{wrMessageId           }  & 2  & 10 & CALIBRATED.               \\ \hline
\multicolumn{8}{|c|}{deltaTx (UInteger64)  }  & 8  & 12 & The deltaTx value (\ref{par:deltaTx}) 
							  of the sending port which shall be stored
							  in otherPortDeltaTx (\ref{par:otherPortDeltaTx})
							  of the receiving port.
% 							  The value of $\Delta_{tx_m}$ measured in 	
% 						   	  picoseconds and multiplied by ${2^{16}}$.
							  \\ \hline
\multicolumn{8}{|c|}{deltaRx (UInteger64) }  & 8  & 20 & The deltaRx value (\ref{par:deltaRx}) 
							  of the sending port which shall be stored
							  in otherPortDeltaRx (\ref{par:otherPortDeltaRx})
							  of the receiving port. 
% 							  The value of $\Delta_{rx_m}$ measured in
% 							  picoseconds and multiplied by ${2^{16}}$.
							  \\ \hline
\end{tabular}
\label{tab:wrCalibratedTLV}
\end{table}




\paragraph{WR\_MODE\_ON} Message sent by the WR Master to the WR Slave. It indicates the successful completion 
of the WR Link Setup process and requests the WR~Slave to validate its WR~mode (by setting wrModeOn 
to TRUE). The message shall have the format specified in Table~\ref{tab:wrOtherTLV}.

\newpage

\subsection{PTP State Machine}
\label{ptpStateMachine}

The WR Link Setup (WR State Machine) is performed in the PTP\_MASTER state of the PTP state machine 
on a WR Port in WR Master mode (WR Master). On a WR Port in WR Slave mode (WR Slave), 
the WR State Machine is executed in the PTP\_UNCALIBRATED state of the PTP state machine. 
Therefore, it is essential for a WR port to implement the PTP\_UNCALIBRATED state as specified in the 
PTP state machine (\textcolor{blue}{1} in Figure~\ref{fig:modifiedPtpFSM}): a transition state 
between the  LISTENING or PRE\_MASTER or MASTER or PASSIVE state and the SLAVE state.
Additionally, WRPTP defines the implementation-specific \textit{SYNCHRONIZATION\_FAULT} and  \\
\textit{MASTER\_CLOCK\_SELECTED} PTP state transition events, see Figure~\ref{fig:modifiedPtpFSM}.


\subsubsection{MASTER\_CLOCK\_SELECTED}
\label{sec:masterClockSelected}

On successful completion of WR State Machine execution (transition from WR\_LINK\_ON to the IDLE 
state with wrModeOn set to TRUE) on the port being in the WR Slave mode and PTP\_UNCALIBRATED state, 
the MASTER\_CLOCK\_SELECTED event (\textcolor{blue}{2} in Figure~\ref{fig:modifiedPtpFSM}, 
defined in Clause 9.2.6.13 of PTP) shall occur. Consequently, the PTP\_SLAVE state shall be 
entered.

\subsubsection{SYNCHRONIZATION\_FAULT}
\label{sec:synchronizationFault}

The port being in the WR Slave mode and PTP\_SLAVE state shall evaluate values of two WR-specific 
parameters: $wrModeOn$ (\ref{par:wrModeON}) and $parentWrModeOn$ (\ref{par:parentWrModeON}). 
If the value of at least one of the parameters is FALSE, the SYNCHRONIZATION\_FAULT event 
(\textcolor{blue}{3} in Figure~\ref{fig:modifiedPtpFSM}, defined in Clause 9.2.6.12 of PTP) shall 
occur and PTP\_UNCALIBRATED shall be entered. Consequently the WR Link Setup is performed. 
Such implementation enables forced re-calibration of WR ports which can be triggered by both, 
WR Master and WR Slave.

 

% \subsubsection{CALIBRATION\_FORCED\_BY\_SLAVE}
% 
% WRPTP introduces a new state transition event to the original PTP state machine: 
% \textit{CALIBRATION\_FORCED\_BY\_SLAVE}. It is depicted in red (number \textcolor{blue}{2})
%  in Figure~\ref{fig:modifiedPtpFSM}. The event shall be instantiated by the reception of
%  the $SLAVE\_PRESENT$ Signaling Message on a WR-Master-enabled port (the value of portDS.wrConfig 
%  is WR\_M\_AND\_S or WR\_M\_ONLY) in $PTP\_MASTER$ state. 
% This transition allows a WR Slave to start WR Master mode and trigger a WR Link Setup process in 
% another WR port.

\newpage

\begin{figure}[ht!]
  \centering
  \includegraphics[width=0.85\textwidth]{protocol/modifiedPtpFSM.ps}
  \caption{Modified PTP Finite State Machine.}
  \label{fig:modifiedPtpFSM}
\end{figure}



\newpage

\subsection{White Rabbit State Machine}
\label{wrFSM}

The White Rabbit Finite State Machine (WR FSM) controls the process of establishing a White Rabbit 
link between a WR Master and a WR Slave (WR Link Setup). It shall be non-preemptive with regards 
to the execution of the PTP State Machine. This means that, once the IDLE state of the WR State Machine is 
exited (WR State Machine execution started), no transitions on the PTP State Machine shall occur 
until the WR State Machine finishes execution and returns to the IDLE state. 

WR Link Setup involves the recognition of two 
compatible WR ports, syntonization over the physical layer, optional measurement of fixed delays and 
exchange of their values across the link. The procedure differs between WR Master and WR Slave. 
Therefore three states of the WR FSM are Slave-only (entered only if the port is in WR Slave mode) 
and one state is Master-only (entered only if the port is in WR Master mode). The WR FSM shall be 
executed in the PTP UNCALIBRATED state on the WR Slave (WR Port in WR Slave mode) and in 
the PTP MASTER state on the WR Master (WR Port in WR Master mode).
The WR FSM is depicted in Figure~\ref{fig:wrFSM} and described in the rest of this section.

Successful completion of WR FSM results in the validation (\textit{wrModeOn} set to TRUE) 
of the recommended WR mode (\textit{wrMode}) and WRPTP synchronization. If EXC\_TIMEOUT\_RETRY 
(section \ref{sec:wrEventsAndConditions}) occurs, it means that WR Link Setup failed, 
WRPTP synchronization cannot be established and standard PTP is performed.

A typical flow of a complete WR Signaling Message exchange between a WR Slave and a WR Master 
during WR Link Setup is presented in Appendix~\ref{ap:wr_lsu_flow}. For clarity, the cooperation of 
PTP and WR state machines from the power up is presented in Appendix~\ref{ap:ptpAndWrFSMS}.

\subsubsection{Conditions to enter WR Slave mode and start the WR FSM}
\label{sec:wrSlaveFSMstart}


%The WR FSM in 
A WR port becomes a WR Slave (portDS.wrMode = WR\_SLAVE) and starts execution 
of the WR FSM by entering the \textit{PRESENT} state, only when the following conditions are met:
\begin{itemize}
\item the port is WR Slave-enabled: \\ 
      \textit{portDS.wrConfig = ($WR\_S\_ONLY$ \textbf{OR} $WR\_M\_AND\_S$)}\\ 
      \textbf{AND}
\item the PTP state machine enters the PTP UNCALIBRATED state as a result of \\
      STATE\_DECISION\_EVENT:
      Recommended State == BMC\_SLAVE (\textcolor{blue}{1} in Figure~\ref{fig:modifiedPtpFSM})
      \textbf{OR} SYNCHRONIZATION\_FAULT event\\
      \textbf{AND}
\item the parent port is WR Master-enabled: \\ 
      \textit{portDS.parentWrConfig = ($WR\_M\_ONLY$ \textbf{OR} $WR\_M\_AND\_S$)} \\ 
      \textbf{AND}
\item at least one of the ports on the link is not in WR Mode:  \\ 
      \textit{portDS.wrMode = FALSE \textbf{OR} portDS.ParentPortWrMode = FALSE}.
\end{itemize}

\newpage

\subsubsection{Conditions to enter WR Master mode and start the WR FSM}
\label{sec:wrMasterFSMstart}

A WR port becomes a WR Master (portDS.wrMode = WR\_Master) and starts execution of WR FSM
by entering the \textit{M\_LOCK} state only when the following conditions are met:
\begin{itemize}
\item the port is WR Master-enabled:\\
      \textit{portDS.wrConfig = $WR\_M\_ONLY$ \textbf{OR} $WR\_M\_AND\_S$} \\ 
      \textbf{AND}
\item the port is in PTP\_MASTER state \\
      \textbf{AND}
\item the $SLAVE\_PRESENT$ Signaling Message has been received. 
% \item the PTP state machine enters the a PTP\_UNCALIBRATED state as a result of 
%       CALIBRATION\_FORCED\_BY\_SLAVE transition event (\textcolor{blue}{2} in 
%       Figure~\ref{fig:modifiedPtpFSM})
\end{itemize}


%The WR FSM in a WR Master-enabled port (\textit{$WR\_M\_ONLY$ \textbf{OR} $WR\_M\_AND\_S$})
%exits the IDLE state and starts execution by entering the \textit{M\_LOCK} state only when 
%PTP state machine enters the PTP\_UNCALIBRATED state as a result of 
%CALIBRATION\_FORCED\_BY\_SLAVE transition event (\textcolor{blue}{2} in 
%igure~\ref{fig:modifiedPtpFSM}).


%\paragraph{State Description}
\subsubsection{State Description}

Table~\ref{tab:wrFSMdesc} specifies the WR states notation used in Figure~\ref{fig:wrFSM}.

\newpage

\begin{table}[hp!]
\caption{WR state definition}
\centering
\begin{tabular}{| c | p{12.2cm} |}          \hline
\textbf{PTP portState}  &  \textbf{Description} \\ 
%&   \\ 
\hline
\small
% IDLE               &  The WR FSM shall be in the IDLE state if the PTP FSM is in a state other than 
% 		      UNCALIBRATED or MASTER. \\ \hline
IDLE               &  The WR FSM shall be in the IDLE state while WR Link Setup is not being performed. 
		      \\ \hline
PRESENT            &  Slave-only state. Upon entering this state, the WR Slave sends a SLAVE\_PRESENT
		      message to the WR Master and waits for the $LOCK$ message.\\ \hline
M\_LOCK            &  Master-only state. Upon entering this state, the WR Master sends the $LOCK$ 
		      message. In this state, the WR Master waits for the WR Slave to finish 
		      successfully the locking process (reception of the $LOCKED$ message). \\ \hline
S\_LOCK            &  Slave-only state. The WR Slave locks its clock signal to the frequency  distributed
		      over the physical layer by the WR Master. \\ \hline
LOCKED             &  Slave-only state. Upon entering this state, the WR Slave sends the $LOCKED$ message 
		      to inform that it is syntonized, and waits for the \textit{CALIBRATE} message.
		      \\ \hline
CALIBRATION 	   &  In this state, optional calibration of the port's reception and/or 
		      transmission fixed delays can be performed. 
		      Upon entering this state, the WR Port sends a \textit{CALIBRATE} message to 
		      the other WR Port. If the calibration of reception fixed delay is needed,
		      the \textit{calSendPattern} flag in the \textit{CALIBRATE} message 
		      is set to TRUE (0x1). If the calibration is not needed, the
		      \textit{calSendPattern} flag is set to FALSE (0x0).  
		      If calibration is not needed, the next state is entered, otherwise an
		      indication from the hardware
		      that the calibration has been finished successfully is awaited. \\ \hline
CALIBRATED         &  Upon entering this state, the WR Port sends a \text{CALIBRATED} message with the 
		      values of its fixed delays. \\ \hline
RESP\_CALIB\_REQ   &  The WR Port's action in this state depends on the value of the 
		      \textit{calSendPattern} flag received in the \textit{CALIBRATE} 
		      message. A TRUE value of the flag
		      indicates that the calibration pattern shall be enabled while the WR State 
		      Machine is in this state. The pattern shall  be
		      disabled on exiting the state (after the timeout of \textit{calPeriod} or 
		      on reception of the \textit{CALIBRATED} message). If the value of  the
		      \textit{calSendPattern} flag is FALSE, 
		      the \textit{CALIBRATED} message is awaited for a timeout of calPeriod or 
		      wrStateTimeout (if calPeriod is 0x0). On 
		      reception of the \textit{CALIBRATED} message, the next state is entered.\\
		      \hline
WR\_LINK\_ON       &  Upon entering this state, the WR Master sends the WR\_LINK\_ON message. 
		      In this state, the values of \textit{wrModeOn (\ref{par:wrModeON})} and 
		      \textit{parentWrModeOn (\ref{par:parentWrModeON})} are set to 
		      TRUE\footnotemark[6] and the \textit{IDLE}  state is entered unconditionally. 
		      The execution of the WR FSM is considered to be completed successfully. \\ \hline
\end{tabular}
\label{tab:wrFSMdesc}
\end{table}

\footnotetext[6]{Depending on the implementation, it might be necessary to update the wrModeOn 
		 field of the best Foreign Master record (clause 9.3.2.4.1 of PTP).}

\addtocounter{footnote}{1}

\begin{figure}[ht!]
  \centering
%  \vspace{-1.3cm}
  \includegraphics[width=1.00\textwidth]{protocol/wrFSM.eps}
  \caption{White Rabbit state machine.}
  \label{fig:wrFSM}
\end{figure}



\newpage


%\paragraph{WR FSM Transition Events and Conditions}
\subsubsection{WR FSM Transition Events and Conditions}
\label{sec:wrEventsAndConditions}
\begin{description}
\item[POWERUP]				Turning on power to the device or resetting.

\item[WR LINK SETUP REQUIRED DECISION] 	(abrv. D\_WR\_SETUP\_REQ) 	Event indicating that a 
      Link Setup is required and that the WR FSM should be executed starting at the
      \textit{PRESENT} state. It occurs when the condition presented in 
      section~\ref{sec:wrSlaveFSMstart} is fulfilled.

\item[SLAVE PRESENT MESSAGE] 		(abrv. M\_SLAVE\_PRESENT) 	Event triggered by the 
      reception of the $SLAVE\_PRESENT$ Signaling message. It requests the recipient to enter 
      the \textit{M\_LOCK} state.

\item[LOCK MESSAGE] 			(abrv. M\_LOCK) 		Event triggered by the 
      reception of the $LOCK$ Signaling message.
% which triggers frequency locking over the physical layer.

\item[LOCKED HARDWARE EVENT] 		(abrv. HW\_LOCKED) 		Event triggered by the 
      reception of the hardware event notification (Table~\ref{tab:outputWrCommWithHW}) indicating 
      that frequency locking has been completed successfully.

\item[LOCKED MESSAGE] 			(abrv. M\_LOCKED)		Event triggered by the 
      reception of the $LOCKED$ Signaling Message which notifies the WR Master that frequency 
      locking has been completed successfully by the WR Slave.

\item[CALIBRATE MESSAGE] 		(abrv. M\_CALIBRATE) 		Event triggered by the 
      reception of the $CALIBRATE$ Signaling message. It requests the recipient to enter the \\
      \textit{RESP\_CALIB\_REQ} state.
% , indicates whether sending of the calibration pattern is
%       required and provides calibration parameters.

% \item[CALIBRATED HARDWARE EVENT] 	(abrv. HW\_CALIBRATED) 		Notification from the
%       hardware indicating that calibration has been completed successfully. 

\item[CALIBRATED MESSAGE] 		(abrv. M\_CALIBRATED)		Event triggered by the 
      reception of the  \text{CALIBRATED} Signaling message. It indicates that the port is 
      calibrated. If the \textit{CALIBRATED} message is received when the calibration pattern is 
      being sent by the recipient, sending of the pattern shall be disabled. 
% The message carries information about fixed delays
%       of the sending port.

\item[WR MODE ON] 			(abrv. M\_WR\_MODE\_ON)		Event triggered by the 
      reception of the \textit{WR\_MODE\_ON} Signaling message. It indicates that the WR Master 
      finished successfully the WR Link Setup and set the \text{wrModeOn} flag to TRUE. 
%       It requests the WR Slave to set the \text{wrModeON} flag to TRUE.

\item[RETRY $n_{name}$]	The White Rabbit state
      machine waits in a given state for a transition event only for a limited time, called timeout.
      The default timeout for the following states is defined by \textit{wrStateTimeout} 
      (\ref{par:wrStateTimeout}): \textit{M\_LOCK}, \textit{CALIBRATED}, \textit{PRESENT}, 
      \textit{S\_LOCK}, \textit{LOCKED}, \textit{CALIBRATION} and \textit{RESP\_CALIB\_REQ}.\\
      The timeout for the  \textit{CALIBRATION} state shall be overwritten by the 
      calPeriod \\ (\ref{par:calPeriod}), if its value is greater than 0x0. Otherwise 
      \textit{wrStateTimeout} shall be used.
      The timeout for the \textit{RESP\_CALIB\_REQ} shall be overwritten by the otherPortCalPeriod 
      (\ref{par:otherPortCalPeriod}), if its value is greater than 0x0. Otherwise 
      \textit{wrStateTimeout} shall be used.\\
      After the timeout expires, the state is re-entered. A state can be re-entered for a maximum 
      number of   
      $n_{\{M\_LOCK, CALIBRATION, CALIBRATED, PRESENT, S\_LOCK, LOCKED, RESP\_CALIB\_REQ\}}$ times.\\
      The default maximum number of state re-entering (retries) is defined by \textit{wrStateRetry} 
      (\ref{par:wrStateRetry}). The number of $n_{CALIB}$ shall be overwritten by \textit{calRetry} 
      (\ref{par:calRetry}) and $n_{RESP\_CALIB\_REQ}$ shall be overwritten by otherPortCalRetry 
      (\ref{par:otherPortCalRetry}) if their values are greater then 0x0. Otherwise, 
      \textit{wrStateRetry} shall be used.

\item[EXCEED TIMEOUT RETRIES] 		(abrv. EXC\_TIMEOUT\_RETRY)	Indicates that a state 
      has been re-entered for a maximum number of times \\
      ($n_{\{M\_LOCK, CALIBRATION, CALIBRATED, PRESENT, S\_LOCK, LOCKED, RESP\_CALIB\_REQ\}}$) and 
      the $n^{th}$ timeout has expired.
									  
\end{description}

\newpage

\subsection{Communication between WR State Machines}
\label{sec:communicationWRfsm}

This section clarifies the communication of the WR State Machines executed on different WR Ports 
on the same link during the WR Link Setup. The communication is performed using WR Signaling Messages
(section~\ref{wrSignalingMSG}). In particular, each WR Signaling Message is sent by a WR Port 
upon entering a particular state of WR State Machine. The same message received on the other WR~Port 
triggers a transition in the WR State Machine (section~\ref{sec:wrEventsAndConditions}).

Table~\ref{tab:wrMsgVsEvents} lists the WR Signaling Messages associating them with the moment 
of their transmission and the event they trigger on the reception WR Port. 
It also specifies the WR Mode of the sending port.


\begin{table}[ph!]
\caption{Communication between WR State Machines on different Boundary Clocks.}
\centering
\begin{tabular}{| l | c | p{2.5cm} | c |}          \hline

\textbf{WR Message} &\textbf{Sent on entering WR state} & \textbf{Sent by port} & 
\textbf{Triggers event}    \\ 
               &             	  & \textbf{in wrMode} 	  &  			           \\ \hline
SLAVE\_PRESENT & PRESENT     	  & WR Slave		  & M\_SLAVE\_PRESENT	           \\ \hline
LOCK           & M\_LOCK      	  & WR Master		  & M\_LOCK		           \\ \hline
LOCKED         & LOCKED	     	  & WR Slave		  & M\_LOCKED		           \\ \hline
CALIBRATE      & REQ\_CALIBRATION & WR Slave or 	  & M\_CALIBRATE	           \\ 
	       & 		  & WR Master 		  & 			           \\ \hline
CALIBRATED     & CALIBRATED	  & WR Slave or 	  & M\_CALIBRATED	           \\ 
	       & 		  & WR Master 		  & 			           \\ \hline
WR\_MODE\_ON   & WR\_LINK\_ON     & WR Master 		  & M\_WR\_MODE\_ON	           \\ \hline
\end{tabular}
\label{tab:wrMsgVsEvents}
\end{table}

\newpage

\subsection{Communication between the WR State Machine and the hardware}
\label{sec:communicationWithHW}

White Rabbit extends the standard PTP communication with the hardware which includes 
transmission/reception of messages and retrieval of timestamps.
The WR communication with the hardware is needed during the WR Link Setup process 
(by the WR State Machine) to control syntonization and fixed delay measurement and to retrieve 
their values. Additionally, outside of the WR Link Setup process (throughout the normal operation 
of the PTP State Machine) a quick detection of the link-down event is necessary. 

\subsubsection{Input to hardware}
\label{sec:inputHW}

The information to the hardware by the WR State Machine is sent on the exiting or entering 
a state of this machine, as described in the Table~\ref{tab:inputWrCommWithHW}. The inputs notify 
the hardware about the need to start/stop a process (i.e. syntonization, pattern transmission).

\subsubsection{Output from hardware}
\label{sec:hwOutput}

The hardware shall communicate two types of information to the WR State Machine
\begin{itemize}
\item event notification, i.e. termination of fixed delays measurement, completion of phase locking or
      disconnection of the cable. 
\item parameter value, i.e. fixed delays. 
\end{itemize}

An event notification shall be evaluated in a particular state of the WR State Machine (defined in 
Table~\ref{tab:outputWrCommWithHW}) by means of polling or interrupt. Its occurrence in other 
state than defined in Table~\ref{tab:outputWrCommWithHW} shall be ignored.

The hardware output of a parameter value (i.e. deltaTx, deltaRx) shall be evaluated on the 
occurrence of an associated event notification (i.e. HW\_CALIBRATED) 
and its value shall be stored in the local Data Set, as defined in Table~\ref{tab:outputWrCommWithHW}.

\newpage

\begin{table}[ht!]
\caption{Input to hardware.}
\centering
\begin{tabular}{| p{5.1cm} | p{2cm} | p{5.5cm} | p{2.6cm} |}          \hline

\textbf{Hardware input}       & \textbf{Type }     & \textbf{Sent on }         & \textbf{Sent by port}  \\ 
                     &                    &         	 	      & \textbf{in wrMode} \\ \hline
\small
HW\_START\_LOCKING   & event notification & entering the 
					    S\_LOCK State          & WR Slave           \\ \hline
HW\_START\_CALIBRATION\footnotemark[7]
		      & event notification & entering the CALIBRATION  
					    State		       & WR Slave and WR Master \\ \hline
HW\_ENABLE\_PATTERN  & event notification & entering the CALIBRATION and 
					    RESP\_CALIB\_REQ States & WR Slave and WR Master \\ \hline
HW\_DISABLE\_PATTERN & event notification & exiting the CALIBRATION and 
					    RESP\_CALIB\_REQ States & WR Slave and WR Master           \\ \hline
\end{tabular}
\label{tab:inputWrCommWithHW}
\end{table}

\footnotetext[7]{Calibration is optional and implementation specific.}
\addtocounter{footnote}{1}


\begin{table}[ht!]
\caption{Output to hardware.}
\centering
\begin{tabular}{| p{5.1cm} | p{2cm} | p{3.5cm} | p{1.7cm} | p{2.4cm} |}          \hline

\textbf{Hardware output}
            &\textbf{Type}& \textbf{Evaluated}    & \textbf{Evaluated by port} & \textbf{Action} \\ 
            &             &         	          & \textbf{in wrMode}         &           \\ \hline
\small
HW\_CALIBRATED\footnotemark[7]
		& event notification    & in the CALIBRATION 
			    State		  & WR~Slave and WR~Master& $calibrated$ Data Set field is set to TRUE, consequently WR State Machine transition is triggered\\ \hline
HW\_LOCKED      & even notification     & in the S\_LOCK State   & WR~Slave& Triggers transition in WR State Machine\\ \hline
HW\_LINK\_DOWN  & even notification     & always   & WR~Slave and WR~Master & Defined in \ref{sec:LinkDown}.\\ \hline
%Set wrModeON to FALSE, and wrMode to NON\_WR \\ \hline
$deltaTx$\footnotemark[7]
		& UInteger64~\footnotemark[8]  & on reception of 
                            HW\_CALIBRATED        & WR~Slave and WR~Master& Save the value in deltaTx DS field.\\ \hline
$deltaRx$\footnotemark[7]
		& UInteger64~\footnotemark[8]  & on reception of 
                            HW\_CALIBRATED        & WR~Slave and WR~Master& Save the value in deltaRx DS field. \\ \hline
% $t_{1}$	    & Integer64   & on sending Sync message& WR~Master             & Send in FollowUp (nanosecond part in correctionField)\\ \hline
% $t_{p2}$    & Integer64   & on receiving Sync message & WR~Slave              & Save\\ \hline
% $t_{3}$     & Integer64   & on sending Delay\_Req message & WR~Slave              & Save\\ \hline
% $t_{p4}$    & Integer64   & on receiving Delay\_Req message & WR~Master             & Send in Delay\_Resp (nanosecond part in correctionField)\\ \hline


\end{tabular}
\label{tab:outputWrCommWithHW}
\end{table}

\footnotetext[8]{The value of $\Delta_{tx_m,rx_s,rx_m,tx_s}$ measured in picoseconds and 
multiplied by $2^{16}$.}
\addtocounter{footnote}{1}

\newpage

\subsection{Link Down}
\label{sec:LinkDown}

Standard PTP can accommodate a temporary connection loss on a link performing the protocol. Therefore 
detection of link down does not mean instant break in the operation of the PTP protocol. Appropriate 
action is taken only after a timeout expires (i.e. PTP State Machine transition and 
Data Sets updates). 
In WR, any loss of connection on a link performing WRPTP results in syntonization failure 
and requires re-establishing the WR Link. Therefore, on reception of the HW\_LINK\_DOWN hardware event 
notification (\ref{sec:hwOutput}), the following actions shall be performed:
\begin{itemize}
  \item wrModeOn (\ref{par:wrModeON}) shall be set to the initialization value (Table~\ref{tab:wrDS}).
  \item wrMode (\ref{par:wrMode}) shall be set to the initialization value (Table~\ref{tab:wrDS}).
  \item calibrated (\ref{par:calibrated}) shall be set to the initialization value (Table~\ref{tab:wrDS}).
  \item the record of Foreign Masters on the port on which the event occurred shall be cleared.
\end{itemize}


\subsection{Re-establishing the WR Link}
\label{sec:reestablishingWRLink}

It shall be possible for an external event (user demand, management message) which occurs on 
a WR Port in an \textit{active} WR Mode (wrModeOn=TRUE, section~\ref{par:wrMode}) 
to enforce re-establishing of the WR Link, i.e. by performing the WR Link Setup. 
 
The WR\_MODE\_OFF event instantiation is WR-implementation-specific. This event should be
instantiated whenever a WR Port is in the IDLE state of the WR State Machine (regardless of the
PTP state) and the implementation circumstances are detected that require re-establishing 
the WR Link.

The WR\_MODE\_OFF event results in setting the wrModeOn DS field of a given WR Port to FALSE.

% \subsection{Master-only WR ordinary clocks}
% \label{sec:masterOnly}
% 
% A WR Node may be designated to be a master-only. A master-only WR Node shall be a WR-compatible 
% non-slave-only (as defined in PTP) ordinary clock with the following changes to a default WR 
% configuration:
% \begin{itemize}
%   \item wrConfig   = WR\_M\_ONLY.
%   \item clockClass = 70 (see \ref{sec:wrptpProfile}).
% \end{itemize}

\newpage

\subsection{White Rabbit PTP Profile Summary}
\subsubsection{Identification}
\label{sec:wrptpProfile}

\begin{table}[ph!]
\caption{Profile print form (clause 19.3.3 of PTP)}
\centering
\begin{tabular}{| c | c |}          					\hline
\multicolumn{2}{|c|}{\textbf{PTP Profile}}  				  \\ \hline
profleName           &  White Rabbit     				  \\ \hline
profileVersion       &  1.0               				  \\ \hline
profileIdentifier    &  08-00-03-00-01-00 				  \\ \hline
organizationName     &  European Organization for Nuclear Research (CERN) \\ \hline
sourceIdentification &  http://www.ohwr.org/projects/white-rabbit         \\ \hline
\end{tabular}
\label{tab:wrPtpProfile}
\end{table}

\subsubsection{PTP attribute values}

All nodes shall support the ranges and shall have the default initialization values for 
the attributes as follows:
\begin{itemize}
  \item portDS.logSyncInterval: The default initialization value shall be 0. The configuration range
	shall be -1 to 6.
  \item defaultDS.priority1: The default initialization value shall be 64.
  \item defaultDS.domainNumber: The default initialization value shall be 0. Only the default 
	domain is allowed.
\end{itemize}

\subsubsection{PTP Options}

All options of 15.5.4.1.7 and clause 17 of PTP are permitted. By default, these options shall be 
inactive unless specifically activated by a management procedure.
\\
The node management shall implement the management message mechanism of the IEEE1588-2008 standard.
\\
The best master algorithm shall be the algorithm specified in section~\ref{sec:wrBMC} of this 
document (Modified BMC).
\\
The delay request-response mechanism shall be the only path delay measurement mechanism.
\\
The TLV mechanism described in section~\ref{sec:wrMSG} shall be supported.
\\

% \subsection{WRPTP Extension to PTP}
% 
% WRPTP extends PTP by adding the state transition event defined in section~\ref{ptpStateMachine}.

\subsection{Implementation-specific additions to PTP required by WRPTP}
 
WRPTP shall implement the following implementation-specific features:
\begin{itemize}
  \item issuing PTP messages with a WR TLV (i.e. suffixed Announce and Signaling messages),
	as described in section~\ref{sec:wrMSG},
  \item proper handling of PTP messages with a WR TLV as described in \ref{sec:wrMSG},
  \item a non-preemptive WR state machine as defined in \ref{wrFSM},
  \item WR data set and additional WR-specific fields to PTP data sets as defined in 
	\ref{sec:wrDS},
  \item SYNCHRONIZATION\_FAULT state transition as defined in \ref{sec:synchronizationFault},
  \item MASTER\_CLOCK\_SELECTED state transition as defined in \ref{sec:masterClockSelected},
  \item communication with hardware as described in \ref{sec:communicationWithHW}.
\end{itemize}


\newpage
\begin{center}
\huge Appendix
\end{center} 

\normalsize
\appendix
\section{PTP State Machine}
\label{ptpFSM}

 \begin{figure}[ht!]
   \centering
   \includegraphics[width=0.8\textwidth]{protocol/ptpFSM.ps}
   \caption{State machine for a full implementation of PTP (Figure 23, IEEE1588).}
   \label{fig:ptpFSM}
 \end{figure}
 
 \newpage



\begin{table}[hp!]
\caption{PTP portState definition (Table 10, IEEE1588).}
\centering
\begin{tabular}{| c | p{9.5cm} |}          \hline
\textbf{PTP portState}  &  \textbf{Description} \\ 
&   \\ \hline
\small
INITIALIZING       &  \small While a port is in the INITIALIZING state, the port initializes its 
		      data sets, hardware, and
		      communication facilities. No port of the clock shall place any PTP messages
		      on its
		      communication path. If one port of a boundary clock is in the INITIALIZING 
		      state, then all ports
		      shall be in the INITIALIZING state.  \\ \hline
FAULTY             &  \small The fault state of the protocol. A port in this state shall not place 
		      any PTP messages except for
		      management messages that are a required response to another management 
		      message on its
		      communication path. In a boundary clock, no activity on a faulty port  
		      shall affect the other ports
		      of the device. If fault activity on a port in this state cannot be confined 
		      to the faulty port, then all
		      ports shall be in the FAULTY state. \\ \hline
DISABLED           &  \small The port shall not place any messages on its communication path. In a
		      boundary clock, no
		      activity at the port shall be allowed to affect the activity at any other
		      port of the boundary clock.
		      A port in this state shall discard all PTP received messages except for 
		      management messages. \\ \hline
LISTENING          &  \small The port is waiting for the announceReceiptTimeout to expire or to
		      receive an Announce
		      message from a master. The purpose of this state is to allow orderly addition 
		      of clocks to a
		      domain. A port in this state shall not place any PTP messages on its 
		      communication path except
		      for Pdelay\_Req, Pdelay\_Resp, Pdelay\_Resp\_Follow\_Up, or signaling 
		      messages, or management
		      messages that are a required response to another management message. \\ \hline
PRE\_MASTER         & \small The port shall behave in all respects as though it were in the 
		      MASTER state except that it shall
		      not place any messages on its communication path except for Pdelay\_Req, 
		      Pdelay\_Resp,
		      Pdelay\_Resp\_Follow\_Up, signaling, or management messages.  \\ \hline
MASTER             &  \small The port is behaving as a master port.   \\ \hline
PASSIVE            &  \small The port shall not place any messages on its communication path except 
		      for Pdelay\_Req,
		      Pdelay\_Resp, Pdelay\_Resp\_Follow\_Up, or signaling messages, or management
		      messages that
		      are a required response to another management message.\\ \hline
UNCALIBRATED       &  \small One or more master ports have been detected in the domain. The 
		      appropriate master port has
		      been selected, and the local port is preparing to synchronize to the selected 
		      master port. This is a
		      transient state to allow initialization of synchronization servos, updating of
		      data sets when a new
		      master port has been selected, and other implementation-specific activity.   
		      \\ \hline
SLAVE              &  \small The port is synchronizing to the selected master port.  \\ \hline
\end{tabular}
\label{tab:ptp_port_state_def}
\end{table}


\newpage

\section{Example White Rabbit Hardware Support Implementation}
\label{sec:sampleHW}

%\subsection{Introduction}

A White Rabbit switch, described in Tomasz W\l{}ostowski's Master Thesis \cite{tomekMSC}, 
implements WRPTP and WR Hardware Support for Gigabit Ethernet over fiber. 
The most important fragments of the document are cited below (with necessary modifications).
The citation is approved by the author.

% Figure \ref{fig:ptp_wr_flow} shows the order of the PTP message exchanges
% during all the phases of the synchronization process, indicating which
% messages are standard PTP (black) and WR-specific (red). The drawing is an overview of 
% the entire WR-PTP synchronization flow described in below.
% 
% \begin{figure}[htb]
%   \centering
%   \includegraphics[width=0.8\textwidth]{fig/tomeksDrawings/ptp_wr_flow.eps}
%   \caption{PTP message flow during WR synchronization.}
%   \label{fig:ptp_wr_flow}
% \end{figure}


%%%%%%%%%%%%%%%%%%%%%%%%%%%%%%%%%%%%%%%5
 \def\us{\char`\_}


\subsection{WR link model}
\label{s:link_model}

Knowledge of the physical model of the links connecting the clocks is a
prerequisite for achieving the required synchronization accuracy. The model
of a WR optical link is depicted in Figure~\ref{fig:link_model}.
\begin{figure}[ht!]
  \centering
  \includegraphics[width=15cm]{../../../figures/protocol/link_model.eps}
  \caption{Model of a WR link (a) and relations between master and slave
  clocks (b)}
  \label{fig:link_model}
\end{figure}
[...]
% \textit{Delta} values shown in Figure~\ref{fig:link_model} express the delays
% introduced by various elements of the system. Big deltas ($\Delta_{txm},
% \Delta_{rxs},  \Delta_{txs},  \Delta_{rxm}$) indicate constant delays
% (e.g. the delays which never change when the link is active), while small
% deltas ($\delta_{ms}, \delta_{sm}$) express the delays which can vary during
% the operation of the link.

WR employs the clock loopback technique to measure the round-trip phase shift
$phase_{MM}$, which is a starting point for determining the precise two-way
delay $delay_{MM}$ and thus the clock offset $offset_{MS}$. The clock signal is
transferred between the master and the slave according to the following scheme:
\begin{enumerate}
\item The reference clock (1) is used to encode the master's transmitter
output and then extracted from the data stream at the slave's receiver. The
recovered clock (2) is a copy of clock (1) delayed by $delay_{MS}$, where
\begin{equation}
    \label{eq:delayms2}
    delay_{MS} = \Delta_{txm} + \delta_{ms} + \Delta_{rxs}
\end{equation}
expresses the master-to-slave latency introduced by the transceivers and
the link.
\item The recovered clock (2) is fed to a clock adjustment unit which shifts
its phase by a programmable value $phase_{S}$ to obtain the phase-compensated
clock (3) -- the final result of WR synchronization.
\item Clock (3) encodes the slave's outgoing data stream and is recovered
at the master side as signal (4).
\item The master measures the phase shift $phase_{MM}$ between its outgoing
(1) and incoming (4) clocks using a phase detector.
\end{enumerate}
The value of $phase_{MM}$ can be expressed as:
\begin{equation}
    \label{eq:phasemm}
    phase_{MM} = \{\Delta + \delta_{ms} + \delta_{sm} + phase_{S}\} \bmod
    T_{ref}
\end{equation}
where $\Delta = \Delta_{txm} + \Delta_{rxs} + \Delta_{txs} + \Delta_{rxm}$
and $T_{ref}$ is the period of Gigabit Ethernet 125 MHz reference clock (8 ns).

The goal of the presented model is to calculate the precise value of
master-to-slave offset $offset_{MS}$ by combining a coarse timestamp-based
round-trip delay \eqref{eq:meanPath} with precise phase measurement
$phase_{MM}.$ Once the offset is computed, the WR slave can phase-shift its
recovered clock (deriving $phase_{S}$ from $offset_{MS}$) to match the phase
of the master clock, completing the synchronization.
\begin{figure}[ht!]
  \centering
  \includegraphics[width=15cm]{../../../figures/protocol/sync_flow.eps}
  \caption{WR synchronization flow (DMTD is explained in Appendix~\ref{s:dmtd})}
  \label{fig:sync_flow}
\end{figure}
Determining the precise offset, however, is not a trivial task. Figure
\ref{fig:sync_flow} shows the steps needed to achieve and maintain
synchronization of a single WR link. The synchronization process can be
split into two parts: \begin{itemize}
\item \textbf{initial synchronization}, which determines the value of
$offset_{MS}$ and fixes the slave's phase shift to compensate the offset.
\item \textbf{phase tracking} which monitors the changes of $phase_{MM}$
over time and updates the phase shifter in the slave to follow these changes
and sustain synchronization.
\end{itemize}

\newpage

\subsection{Link detection and syntonization}
During the first two steps of the synchronization flow
\ref{fig:sync_flow}, a syntonized WR connection is established. Drawing
\ref{fig:link_detect_and_syntonization} shows the order of hardware operations
and message exchanges which result in a syntonized link.
\begin{figure}[ht!]
  \centering
  \includegraphics[width=7cm]{../../../figures/protocol/link_detect_and_syntonization.eps}
  \caption{WR Link detection and syntonization}
  \label{fig:link_detect_and_syntonization}
\end{figure}
At the initial moment, the master and the slave are not connected to
each other. Their PHYs are transmitting an Ethernet idle pattern (see Table 3.1 \cite{tomekMSC}), 
but are not receiving any meaningful bitstreams. As soon
as the physical connection is present, the PHYs will start receiving valid
idle patterns and the Ethernet link will become active (see \cite{IEEE802.3}
section 36.2.5.2.6). Presence of a valid physical link will trigger the
following sequence:
\begin{enumerate}
\item The master starts broadcasting \textit{ANNOUNCE} messages to look for
a WR slave,
\item Eventually, the slave will respond with a \textit{SLAVE\_PRESENT}
message, indicating that it supports WR. If no response has been received
within a predefined time, the master assumes that the slave is not WR
compatible and aborts the synchronization process,
\item The master issues a \textit{LOCK} message, commanding the slave to
begin recovering the clock from its received data stream,
\item The slave sets up its PLL to use the RX clock as a reference and as
soon as the PLL is locked, responds with a \textit{LOCKED} message.
\end{enumerate}
Upon successful completion of the above process, both nodes are syntonized
and the delay measurements can begin.

\subsection{Coarse delay measurement}
\label{s:coarse_delay}
The coarse delay measurement produces a set of timestamps $t_{1} ... t_{4}$
which are obtained using the PTP End-to-End, two-way clock method as shown
in fig. \ref{fig:wrPTPmsgs}. The timestamps are generated in the hardware
with a circuit depicted in fig. \ref{fig:coarse_measurement}, guaranteeing
single-cycle timestamping accuracy, which is necessary for reliably merging the
coarse delay with the $phase_{MM}$ phase shift during fine delay calculation.
\begin{figure}[ht!]
  \centering
  \includegraphics[width=17cm]{../../../figures/protocol/coarse_measurement.eps}
  \caption{Structure of a WR timestamping unit (TSU)}
  \label{fig:coarse_measurement}
\end{figure}
Timestamps in WR are taken when the PCS detects a Start-of-Frame Delimiter
(SFD) character in the incoming (RX timestamps) or outgoing (TX timestamps)
data stream (see section 3.2.2 \cite{tomekMSC}).

Let's first focus on the blocks marked blue in
Figure~\ref{fig:coarse_measurement}. Each time an SFD character is detected,
the PCS produces a timestamp trigger pulse which causes the timestaming unit
to take a snapshot of the free running counter CNTR\_R with the D-type
register DREG\_R. The counter is counting from 0 to 124999999 which (given
the reference clock frequency of 125 MHz) gives a period of one full second.

The counter CNTR\_R works synchronously to the reference clock (master side,
signal 1 in Figure~\ref{fig:link_model}) or the compensated clock (slave side,
signal 5). Since RX trigger pulses come from different clock domains (2 or
4), they need to be synchronized to the reference clock with a chain of
D flip-flops (SYNC). Single-cycle long trigger pulses are widened by the
pulse extender before going to the synchronizer chain to ensure that no
pulses are missed due to the metastability of synchronizer flip-flops. The
counter value is latched in register DREG\_R on the rising edge of the
synchronizer output. TX timestamp triggers, which are generated within the
reference clock domain (clocks 1 or 5) also pass through a synchronizer
chain to obtain identical trigger reaction latency.

Unfortunately, due to the jitter of clock signals, crossing clock domains
can make the gathered timestamps useless by causing a random $\pm 1$ LSB
error when the RX clock and the reference clocks are almost in phase. The
problem is illustrated in Figure~\ref{fig:ts_jitter}.  Dashed lines show
the transitions of ideal (jitter-free) clocks. If the jitter is neglected,
the reference clock transition should be slightly ahead of the RX clock
transition and the gathered timestamp should be equal to 2. However, if the
clocks are jittery, the transitions may sometimes occur in reverse order,
producing an erroneous timestamp of value 1.
\begin{figure}[ht!]
  \centering
  \includegraphics[width=10cm]{../../../figures/protocol/ts_jitter.eps}
  \caption{Timestamping errors caused by clock jitter}
  \label{fig:ts_jitter}
\end{figure}
One possible way of addressing this issue is to take RX timestamps on both
reference clock edges. The falling edge part of the TSU is marked pink in
\ref{fig:coarse_measurement}. It does not have an independent counter --
instead, the current value of the rising edge counter is latched in register
CNTR\_F on the falling edge of the reference clock, making the CNTR\us
F a copy of CNTR\_R delayed by a half of the clock period. This method
ensures that at least one of the timestamps is valid at any moment 
(see Figure~\ref{fig:ts_dualedge}). The correct timestamp is chosen depending on the
current phase shift between clocks (see section \ref{s:fine_delay}).
\begin{figure}[ht!]
  \centering
  \includegraphics[width=12cm]{../../../figures/protocol/ts_dualedge.eps}
  \caption{Dual-edge timestamping in WR}
  \label{fig:ts_dualedge}
\end{figure}
In order to simplify the hardware design, the presented TSU can only measure
the sub-second part of PTP timestamps. The UTC part is appended in the
software using the algorithm shown in listing \ref{lst:utc_gen}.
\begin{lstlisting}[caption=Producing UTC timestamps,label=lst:utc_gen]
  timestamp tstamp = get_hardware_timestamp();
  time t = get_utc_time();
// check if there was a transition of the UTC counter between the
// generation of the timestamp and its readout by the software
  if(t.milliseconds < tstamp.rising_edge)
   tstamp.seconds = t.seconds - 1;
  else
   tstamp.seconds = t.seconds;
\end{lstlisting}


\subsection{Digital DMTD phase detector}
\label{s:dmtd}
The fine delay measurements in WR are based on a Dual Mixer Time Difference
(DMTD) phase detector. Therefore, in this section a short introduction to
DMTD technology will be presented.
\begin{figure}[ht!]
  \centering
  \includegraphics[width=11cm]{../../../figures/misc/analog_dmtd.eps}
  \caption{Structure of an analog DMTD phase detector}
  \label{fig:analog_dmtd}
\end{figure}
Figure \ref{fig:analog_dmtd} shows an analog DMTD system \cite{allan90}. Let's
assume that its input clocks $a(t)$ and $b(t)$ have identical amplitudes and
frequencies equal to $f_{clk}$. The phases of both clocks are respectively
$\phi_{a}$ and $\phi_{b}$. The clocks are multiplied by the mixers with
the signal $c(t)$ of frequency $f_{offset}$ and phase $\phi_{offset}$. The
multiplication result can be expressed as:

\begin{eqnarray}
    a(t) \cdot c(t) & = & cos(2\pi t f_{clk} + \phi_{a}) \cdot cos (2\pi t
    f_{offset} + \phi_{offset}) \nonumber \\
\label{eq:mixing1}& = & \frac{1}{2} cos(2\pi t (f_{clk} + f_{offset}) +
\phi_a + \phi_{offset}) \\
\label{eq:mixing2}& + & \frac{1}{2} cos (2 \pi t (f_{clk} - f_{offset}) +
\phi_a - \phi_{offset})
\end{eqnarray}

Term \ref{eq:mixing1} has higher frequency than the multiplied signals and
it's removed by the low-pass filter, leaving only the low-frequency term
\ref{eq:mixing2}. The downconversion process changes the frequency of the
signals, but does not change their phase difference. Therefore, if the offset
frequency $f_{offset}$ is close to the input signal frequency $f_{clk}$,
the phase shift can be measured by counting the time between the rising
edges of the downconverted clocks. For example, a reference clock of 125
MHz and an offset clock of 124.99 MHz will produce an output signal of 10
kHz. At such frequencies, the phase shift can be very accurately measured
using a simple counter.

Analog DMTDs provide excellent resolution and linearity, at the cost of several
external discrete components (mixers and filters), which can be troublesome,
especially in multi-port applications such as the WR switch. Fortunately, the
analog mixing operation can be transformed into a digital sampling operation,
resulting in a digital DMTD detector, shown on fig. \ref{fig:digital_dmtd}.
\begin{figure}[ht!]
  \centering
  \includegraphics[width=12cm]{../../../figures/misc/dmtd.eps}
  \caption{Structure of a digital DMTD phase detector}
  \label{fig:digital_dmtd}
\end{figure}
In a Digital DMTD (DDMTD) \cite{icalepcs09}, the input signals are square
waves, the mixers are replaced with simple D-type flip-flops and the offset
clock is generated with a PLL from one of the input clocks. The sampling
operation performed by the flip-flops can be mathematically proved to be
equivalent to analog mixing (\ref{eq:mixing2}), but the principle of a DDMTD
can be explained in a more intuitive way.
\begin{figure}[ht!]
  \centering
  \includegraphics[width=15cm]{../../../figures/misc/dmtd_vernier.eps}
  \caption{A vernier (a) and signals generated by DDMTD (b)}
  \label{fig:dmtd_vernier}
\end{figure}
Figure \ref{fig:dmtd_vernier}a shows a vernier caliper. It has two scales -
the big millimeter scale and a small vernier scale, with the units slightly
smaller than the main scale. For a typical caliper, the vernier scale is split
into 10 intervals with 4.9 mm spacing. If the length of the measured object
has a fractional part, one of bars on the vernier scale will be in line with
one of bars from the main scale. Translating to the language of electronics,
the main scale is the input clock where each interval represents one cycle,
the vernier scale is the offset clock and the transitions in the DDMTD output
signals occur when bars on both scales are aligned. An example of signals
produced by a DDMTD with $N = 8$ is shown in fig. \ref{fig:dmtd_vernier}b. The
output phase shift is proportional to the input phase shift $\phi$ by a
factor of $(N+1)$. A general relation between the phase shift and the cycle
interval between the DMTD outputs can be expressed as:

\begin{equation}
\label{eq:dmtd1}
\phi \mbox{[ns]} = \frac{n_{cycles}}{N+1} \cdot \frac{1}{f_{in}}
\end{equation}


Digital DMTDs have all the advantages of their analog counterparts while
requiring only one external component (the oscillator producing the offset
clock) which is shared among all measurement channels.	This opens the way
for low-cost FPGA implementations. If a picosecond-level accuracy is not
required, even a PLL integrated inside an FPGA can be used, eliminating all
external components.
\begin{figure}[ht!]
  \centering
  \includegraphics[width=8cm]{../../../figures/misc/dmtd_glitches.eps}
  \caption{Glitches in the DMTD output caused by clock jitter}
  \label{fig:dmtd_glitches}
\end{figure}
In practical DDMTD implementations, the output signals need to be additionally
conditioned as the input clock jitter can introduce glitches, as shown
on fig. \ref{fig:dmtd_glitches}. More details about the deglitching and
postprocessing algorithm can be found in section 4.3.5 of \cite{tomekMSC}.

\subsection{Fine delay measurement}
\label{s:fine_delay}
Having the knowledge of PTP timestamps and round-trip phase shift $phase_{MM}$,
we can calculate the precise round-trip delay $delay_{MM}$. The calculation
is performed in two steps:
\begin{itemize}
\item The accuracy of PTP timestamps $t_{2}$ and $t_{4}$ is enhanced beyond a
single clock cycle using the knowledge of the round-trip phase $phase_{MM}$
and the slave PLL setpoint $phase_{S}$. As a result, precise timestamps
$t_{2p}$ and $t_{4p}$ are obtained.
\item The fine round-trip delay is calculated using the standard PTP formula:
\begin{equation}
\label{eq:ptp_precise}
delay_{MM} = (t_{4p} - t_{1}) - (t_{3} - t_{2p})
\end{equation}
\end{itemize}

Only the reception timestamps need to be enhanced, as they are generated
within clock domains asynchronous to the reference (or compensated) clock
(see table \ref{tab:ts_domains}). Transmission timestamps are always integer
because packets are transmitted and timestamped using the same clock.

The flow graph of the algorithm used to merge PTP timestamps with
phase measurements is shown on Figure~\ref{fig:merging_timestamps}. 
Figure~\ref{fig:merging_example} depicts sample measurements of inputs and results
of the enhancement algorithm for $t_{4p}$ timestamps, where the varying link
delay was simulated using the slave's phase shifter.

\begin{table}[htb]
  \caption{Timestamping clock domains}
  \centering
  \begin{tabular}{|c|c|c|c|c|} \hline
    Timestamp & Origin & Trigger clock & Timestamping clock & Correction \\
    \hline \hline
    $t_{1}$ & Master TX & reference (1) & reference (1) & 0 \\ \hline
    $t_{2p}$ & Slave RX & slave recovered (2) & compensated (3) & $-phase_{S}$
    \\ \hline
    $t_{3}$ & Slave TX & compensated (3) & compensated (3) & 0 \\ \hline
    $t_{4p}$ & Master RX & master recovered (4) & reference (1) & $+phase_{S}$
    \\ \hline
  \end{tabular}
  \label{tab:ts_domains}
\end{table}
\begin{figure}[ht!]
  \centering
  \includegraphics[width=13cm]{../../../figures/protocol/merging_timestamps.eps}
  \caption{Algorithm for enhancing coarse timestamps with DMTD phase.}
  \label{fig:merging_timestamps}
\end{figure}
The first step of the algorithm addresses the problem of glitches in
RX timestamps, by choosing either the rising or falling edge timestamp
depending on the phase between the RX clock and the reference clock. The
key parameter of this step, $\phi_{trans}$ provides an approximate
value of $phase_{MM}$ at which a transition should occur in the value of
$t_{4r}$. In Figure~\ref{fig:merging_example}, it is equal to 6.6 ns - it's
the approximate intersection point of $phase_{MM}$ (blue, sawtooth-like
trace) and the transition of the rising edge timestamp (red trace). The
value of $\phi_{trans}$ is a device-specific constant. It can be determined
once during the factory calibration of a WR device or measured upon every
startup by sweeping a full clock period using the built-in phase shifter
and searching for a transition in the RX timestamp value.

If the actual value of $phase_{MM}$ lies within a $\pm 25\%  T_{ref}$
range from the transition point $\phi_{trans}$ (green zone in
\ref{fig:merging_example}), the algorithm will use the $t_{4f}$ timestamp,
otherwise the $t_{4r}$ timestamp will be taken (red zone). Note that because
$phase_{MM}$ is bounded to $\left[0, T_{ref}\right)$, the range checks must
be aware of the jump in the phase between $T_{ref}$ and $0$.

The second step checks if the $phase_{MM}$ value is ahead of the transition
point $\phi_{trans}$, and eventually increases the chosen timestamp by a full
cycle ($T_{ref}$). That ensures the transition in $t_{4}$ will always occur
when $phase_{MM} = \phi_{trans}$, eliminating the risk of transition glitches.

The last part of the algorithm simply adds the picosecond part (which is
the DMTD phase corrected with the transition offset $\phi_{trans}$) to
the coarse deglitched timestamp $t_{4}$. If the resulting picosecond part
($\phi$) is negative, an additional full cycle is added to the result. The
final output of the merging algorithm is shown in \ref{fig:merging_example}
as the thick navy trace.
 \begin{figure}[ht!]
  \centering
  \includegraphics[width=12cm]{../../../figures/protocol/merging_example.eps}
  \caption{Example of $t_{4p}$ timestamp enhancing.}
  \label{fig:merging_example}
\end{figure}
The enhancement operation for the $t_{2}$ timestamp is done in a very
similar way by replacing $t_{4}$ with $t_{2}$ and $phase_{MM}$ with
$phase_{S}$. Note that changing the slave's phase shift $phase_{S}$ will
result in a change of the values of both $t_{2p}$ and $t_{4p}$ (see table
\ref{tab:ts_domains}). Increasing $phase_{S}$ by a certain $\delta$ will
cause $t_{4p}$ to also increase by $\delta$ (assuming that the link delay
stays constant). Simultaneously, the value of $t_{2p}$ will be decreased by
the same amount, as for the slave's RX timestamps, the timestamping clock is
$phase_{S}$ ahead of the trigger clock (so if the phase shift is increased,
the timestamp value goes down). Therefore, the calculated fine delay is not
affected by the changes of $phase_{S}$:

\begin{equation}
\label{eq:ptp_precise2}
delay_{MM} = (t_{4p} - t_{1}) - (t_{3} - t_{2p}) = t_{4p} \cancel{+ phase_{S}}
- t_{1} - t_{3} + t_{2p} \cancel{-phase_{S}}
\end{equation}


\subsection{Link asymmetry estimation}
\label{s:asymmetry}
In order to compute an accurate master-slave offset, one must determine the
asymmetry of the link delay. This asymmetry cannot be measured directly (as
a difference between the overall M-S and S-M delays), as it would require
the nodes to be synchronized prior to the measurement. Therefore it is
only possible to estimate the asymmetry from round-trip delay $delay_{MM}$,
using the knowledge of the properties of the components constituting the link.

In the WR optical link model, the following sources of asymmetry were taken
into consideration:
\begin {itemize}
\item Propagation delays of electronic components and PCB traces (circuit
asymmetry),
\item Asymmetry of optical transceivers (SFPs),
\item Difference between TX and RX wavelengths in the fiber,
\item Internal structure and clocking of the PHY (SerDes) chips.
\end{itemize}
\begin{figure}[ht!]
  \centering
  \includegraphics[width=\textwidth]{../../../figures/protocol/asymmetries.eps}
  \caption{Delay asymmetries in WR optical link.}
  \label{fig:asymmetries}
\end{figure}
Figure \ref{fig:asymmetries} shows the reference asymmetry model used in WR
devices. The device's asymmetric delays (as shown in Figure~\ref{fig:link_model})
are expressed as sums of circuit, SFP and PHY delays between the phase detector
inputs (the PD measuring $delta_{MM}$ on the master side and the PD in the
phase shifting PLL on the slave side) and the SFP optical input/output. The
fiber asymmetry is compensated separately in the slave's PTP servo.
\begin{eqnarray}
\Delta_{tx(m/s)} = \delta_{TX\us PHY(M/S)} + \delta_{TX\us CIR(M/S)} +
\delta_{TX\us SFP(M/S)}
\\
\Delta_{rx(m/s)} = \delta_{RX\us PHY(M/S)} + \delta_{RX\us CIR(M/S)} +
\delta_{RX\us SFP(M/S)} \nonumber
\label{eq:asym_delay}
\end{eqnarray}\\

\subsubsection{Circuit and SFP asymmetry}
Circuit asymmetry stems from differing PCB trace lengths, propagation delays
of the electronic components of the clock distribution system and FPGA logic
cell placement and routing delays. SFP asymmetry is a result of different
reception and transmit delays between the electrical and optical ports of
the SFP transceiver.

These asymmetries can vary with changes of the device's operating conditions
(i.e. temperature, supply voltage). Depending on the synchronization accuracy
requirements, they can be:
\begin{itemize}
\item Treated as time-invariant and measured once during factory setup.
\item Actively compensated, following the changes in temperature and/or supply
voltage read from built-in sensors (requires a model of delay vs. temperature
and voltage),
\item Eliminated by building a system where both master and slave have
identical (or zero) asymmetry and operate in similar conditions.
\end{itemize}

The last method is particularly useful for compensating the delays inside
the FPGA on the way between the PHY (oscillator) and the phase detector,
as they cannot be externally measured during the factory calibration
process. They are also significantly affected by temperature and voltage
changes. Therefore, the only practical way of dealing with the FPGA part
of the circuit asymmetry is to equalize the delays on all phase detector
inputs. This can be done by constraining the routing delays or manually
placing and routing the phase detector block. This approach also allows for
reducing the impact of the temperature-voltage changes on the asymmetry to
a negligible value, as two equal paths placed next to each other would have
very similar temperature/voltage delay coefficients.

SFP asymmetries can be compensated in a similar manner by choosing pairs
of SFP transceivers whose TX and RX delays differ by similar values
(i.e. $\delta_{TX\us SFPM} - \delta_{RX\us SFPM} = \delta_{TX\us SFPS} -
\delta_{RX\us SFPS}$). This can be done in a laboratory system where all
other asymmetric delays are already compensated by selecting SFPs that
together give minimum clock offset.

\subsubsection{WDM fiber asymmetry}
\label{s:wdm_asymmetry}
As mentioned in section 3.2 of \cite{tomekMSC}, WR uses different wavelengths
for transmitting and receiving the data (for example 1550/1310 nm). Due
to chromatic dispersion of the fiber, the refractive indexes for these
wavelengths are slightly different, causing different propagation speeds
(and thus, different delays) between the master and the slave.

The refractive index at a given wavelength can be derived using Sellmeier's
equation  \cite{saleh07} (\ref{eq:sellmeier}):

\begin{equation}
\label{eq:sellmeier}
n^2(\lambda) = 1 + \frac{B_{1} \lambda^2}{\lambda^2 - C_{1}} + \frac{B_{2}
\lambda^2}{\lambda^2 - C_{2}} + \frac{B_{3} \lambda^2}{\lambda^2 - C_{3}}
\end{equation}

where $B_{i}$ and $C_{i}$ are material-specific coefficients.

For a standard G.652 telecom fiber, the refractive indexes are respectively:
$n_{1550} = 1.467$ and $n_{1310} = 1.466$. In order to simplify the asymmetry
calculations in the hardware, the WR specification defines
a custom fiber asymmetry coefficient \eqref{eq:fiber_alpha1} expressing the
ratio between the M-S and S-M fiber propagation delays:

\begin{equation}
\label{eq:fiber_alpha1}
\alpha = \frac{\delta_{ms}}{\delta_{sm}} - 1 = \frac{n_{1550}}{n_{1310}} - 1
\end{equation}
Unfortunately, refractive indexes may vary slightly between different fiber
manufacturers, making the direct calculation of $\alpha$ unreliable. WR
can solve this problem by characterizing the fiber asymmetry using
laboratory measurements of $delay_{MM}$ (done by PTP) and clock offset
$offset_{MS}$ (done using an oscilloscope). The measurements are performed
when all the other asymmetries are already compensated. The value of
$\alpha$ \ref{eq:fiber_alpha3} is calculated by solving equation system
\ref{eq:fiber_alpha2}:
\begin{equation}
\label{eq:fiber_alpha2}
\begin{cases}
delay_{MM} = \delta_{ms} + \delta_{sm} + \Delta \\
offset_{MS} = \frac{\delta_{ms} - \delta_{sm}}{2}
\end{cases}
\end{equation}
where $\Delta$ accounts for all fixed delays in the path (i.e. $\Delta_{txm}
+ \Delta_{rxs} + \Delta_{txs} + \Delta_{rxm}$). Substitution of $\delta_{ms}$
and $\delta_{sm}$ with the results of \ref{eq:fiber_alpha2} gives the final
value of $\alpha$:
\begin{equation}
\label{eq:fiber_alpha3}
\alpha = \frac{delay_{MM}  - \Delta + 2\cdot offset_{MS}}{delay_{MM} -
\Delta - 2\cdot offset_{MS}}
\end{equation}

\subsubsection{Transceiver asymmetry}
\label{s:xcvr_asymmetry}
\label{sec:calibForGigbitE}
\textit{[This chapter has been modified to add clarity]}


Transceiver asymmetry is a result of the internal structure of the SerDes's
circuitry. Most of the SerDes chips nowadays are optimized for low power
consumption and fast locking to the incoming data stream. Unfortunately,
because of these optimizations, PHYs may not keep constant transmit/receive
latencies. The problem is illustrated in Figure~\ref{fig:phy_asymmetry}.
\begin{figure}[ht!]
  \centering
  \includegraphics[width=11cm]{../../../figures/protocol/phy_asymmetry.eps}
  \caption{Random delays in gigabit SerDes devices (a) and blocks causing them
  (b).}
  \label{fig:phy_asymmetry}
\end{figure}
PHY asymmetry manifests itself as a random latency between the rising edge of
the TX/RX clock and the inter-symbol boundary in the transmitted (received)
serial data stream. In most PHYs, this latency can be different for each
PLL/CDR lock cycle, but once the PHY is locked, the delays shall remain
constant. PHYs whose TX/RX delays can change when the link is active are
unsuitable for WR devices.

In the PHYs which have been evaluated for usability in WR hardware (TLK1221
from Texas Instruments and Virtex-6/Spartan-6 GTP transceivers), the following
random delays were identified:
\begin{itemize}
\item RX alignment latency, observed in PHYs which do not correct the RX
clock phase when aligning to the inter-symbol boundary (red, observed for
Xilinx GTP),
\item RX latency resulting from the internal structure of the digital
oversampling CDR (green, observed for TLK1221),
\item TX latency caused by the clock divider between the internal PLL and
the parallel-to-serial register (blue, observed for TLK1221).
\end{itemize}
Alignment latency can be measured every time the link goes up by disabling
automatic comma alignment and bit-shifting the unaligned output data
until a valid 8B10B code sequence is detected (\textit{bitslip trick},
see \cite{Peek2010}). Compensation of the latter two latencies, however,
requires additional calibration logic. An example method (used in the WR
switch) is shown in Figure~\ref{fig:phy_latency_measurement}.
\begin{figure}[ht!]
  \centering
%  \includegraphics[width=\textwidth]{fig/tomeksDrawings/phy_latency_measurement.eps}
  \includegraphics[width=0.60\textwidth]{../../../figures/misc/calibration_1.ps}
  \caption{PHY latency measurement using calibration patterns.}
  \label{fig:phy_latency_measurement}
\end{figure}
Such methods can by applied to the
PHYs whose maximum ($\delta_{PHY(M/S)_{max}}$) and minimum ($\delta_{PHY(M/S)_{min}}$) delays are known 
(provided in the datasheet) and the delay's variation is
\begin{equation}
  \label{eq:fixedDelayVariation}
  \delta_{\{TX, RX\}\_PHY(M/S)_{variable}} \in \langle 0 :\delta_{PHY(M/S)_{max}} - \delta_{PHY(M/S)_{min}} \rangle
\end{equation}
e.g. below 10 bit times in the case of Gigabit Ethernet. 
Therefore, a fixed delay can be expressed as a sum of a constant value 
($\delta_{\{TX, RX\}\_PHY(M/S)_{min}}$) and a variable part ($\delta_{\{TX, RX\}\_PHY(M/S)_{variable}}$) which needs to be measured:
\begin{equation}
  \label{eq:fixedDelay}
  \delta_{\{TX, RX\}\_PHY(M/S)} = \delta_{\{TX, RX\}\_PHY(M/S)_{min}} + \delta_{\{TX, RX\}\_PHY(M/S)_{variable}}
\end{equation}
 


The PHY transmit path is fed with a sequence of RD+ K28.5 characters (1111100000),
effectively producing a 125 MHz square wave on the serial outputs. 
The repeated pattern of five "0" and five "1" is defined by the IEEE 802.3 standard \cite{IEEE802.3} 
as \textit{Low-frequency test pattern} (Appendix~36A.2). The phase
shift between the TX clock and the output bitstream can be measured using
a DDMTD phase detector, giving the value of the TX variable latency 
($\delta_{TX\_PHY(M/S)_{variable}}$). The same method
can be used to measure the RX variable latency ($\delta_{RX\_PHY(M/S)_{variable}}$) 
by commanding the link partner to send
the calibration pattern. Note that since the K28.5 character contains a comma
(5 consecutive ones), a burst of subsequent K28.5 symbols will cause improper
operation of the PHY's comma alignment unit. Therefore, comma alignment must
be disabled during the calibration process.

\subsection{Establishing and maintaining synchronization}
\label{s:estab_sync}
Having obtained the values of round-trip delay $delay_{MM}$ and link asymmetry,
we can calculate the one-way master to slave delay $delay_{MS}$ by solving
equation (\ref{eq:delaymm_full}):

\begin{equation}
delay_{MM} = \Delta + \delta_{ms} + \delta_{sm}
\label{eq:delaymm_full}
\end{equation}

where $\Delta = \Delta_{txm} + \Delta_{rxs} + \Delta_{txs} +
\Delta_{rxm}$. Substituting $\delta_{sm}$ with \ref{eq:fiber_alpha1}
and solving for $\delta_{ms}$, we obtain the one-way fiber delay
\ref{eq:delaymm_full_2}:

\begin{equation}
\delta_{ms} = \frac{1+\alpha}{2+\alpha} (delay_{MM} - \Delta)
\label{eq:delaymm_full_2}
\end{equation}

which after adding the circuit, SFP and PHY asymmetric delays present on the
master to slave path gives the final master-slave delay \ref{eq:delaymm_full_3}
and offset \ref{eq:offset_ms}:

\begin{eqnarray}
\label{eq:delaymm_full_3}delay_{MS} & = & \frac{1+\alpha}{2+\alpha} (delay_{MM}
- \Delta) + \Delta_{txm} + \Delta_{rxs} \\
\label{eq:offset_ms} offset_{MS} & = & t_{1} - t_{2p} - delay_{MS}
\end{eqnarray}
The value of $offset_{MS}$ is the input for the slave's offset adjustment
algorithm which controls the slave's clock servo. The flow diagram of the
algorithm is shown in fig. \ref{fig:adjustment_and_servo}a and an example
servo design can be found in fig. \ref{fig:adjustment_and_servo}b. The
algorithm assumes that the frequency has been already syntonized by means
of Sync-E and only the clock offset needs to be corrected. Offset correction
is split into 3 steps:
\begin{figure}[ht!]
  \centering
  \includegraphics[width=\textwidth]{../../../figures/protocol/adjustment_and_servo.eps}
  \caption{WR slave offset adjustment (a) and clock servo (b)}
  \label{fig:adjustment_and_servo}
\end{figure}
\begin{enumerate}
\item \textbf{UTC time adjustment}: the UTC counter in the servo is
increased (or decreased) by the number of full seconds \ref{eq:corr_utc}
in $offset_{MS}$:
\begin{equation}
corr_{utc} = \lfloor{\frac{offset_{MS}}{1 \mathrm{~s}}}\rfloor
\label{eq:corr_utc}
\end{equation}
\item \textbf{Clock cycle counter adjustment}: The reference clock cycle
counter, which produces the PPS signal is adjusted by the number of full
$T_{ref}$ (8 ns) cycles \ref{eq:corr_cnt}.
\begin{equation}
corr_{cnt} = \lfloor \frac{offset_{MS} - corr_{UTC}}{T_{ref}}\rfloor
\label{eq:corr_cnt}
\end{equation}
\item \textbf{Phase adjustment}: The slave's phase shifter is adjusted with
the remaining sub-cycle part of the offset \ref{eq:corr_phase}
\begin{equation}
corr_{phase} = offset_{MS} - [offset_{MS}]
\label{eq:corr_phase}
\end{equation}
\end{enumerate}
\emph{Voil\`{a}!} Now the slave's clock and PPS signals are synchronized to the
master with sub-nanosecond accuracy. Since the offset can vary with operating
conditions, it is measured at regular intervals and the difference between
subsequent measurements is added to the slave's phase shift to compensate for
phase drift:
\begin{equation}
corr_{phase} = offset_{MS} - offset_{MS\us previous}
\label{eq:corr_phase2}
\end{equation}
Because the phase drift is mainly caused by temperature variations, the rate
of subsequent adjustments can be very low, even in the scale of a single
adjustment per hour.

Note that the way $corr_{phase}$ is calculated requires the phase shifter
to be able to change the phase relatively by any arbitrary value with
no "jumps" in the signal when the value of $corr_{phase}$ crosses the
inter-cycle boundary. For example, one can use a PLL with a phase detector
capable of handling wrap-around phase transitions, but not a programmable
delay line. An example design of such phase detector and PLL is described
in section 4.3.7 of \cite{tomekMSC}.


%  \def\us{\char`\_}


\subsection{WR link model}
\label{s:link_model}

Knowledge of the physical model of the links connecting the clocks is a
prerequisite for achieving the required synchronization accuracy. The model
of a WR optical link is depicted in Figure~\ref{fig:link_model}.
\begin{figure}[ht!]
  \centering
  \includegraphics[width=15cm]{../../../figures/protocol/link_model.eps}
  \caption{Model of a WR link (a) and relations between master and slave
  clocks (b)}
  \label{fig:link_model}
\end{figure}
[...]
% \textit{Delta} values shown in Figure~\ref{fig:link_model} express the delays
% introduced by various elements of the system. Big deltas ($\Delta_{txm},
% \Delta_{rxs},  \Delta_{txs},  \Delta_{rxm}$) indicate constant delays
% (e.g. the delays which never change when the link is active), while small
% deltas ($\delta_{ms}, \delta_{sm}$) express the delays which can vary during
% the operation of the link.

WR employs the clock loopback technique to measure the round-trip phase shift
$phase_{MM}$, which is a starting point for determining the precise two-way
delay $delay_{MM}$ and thus the clock offset $offset_{MS}$. The clock signal is
transferred between the master and the slave according to the following scheme:
\begin{enumerate}
\item The reference clock (1) is used to encode the master's transmitter
output and then extracted from the data stream at the slave's receiver. The
recovered clock (2) is a copy of clock (1) delayed by $delay_{MS}$, where
\begin{equation}
    \label{eq:delayms2}
    delay_{MS} = \Delta_{txm} + \delta_{ms} + \Delta_{rxs}
\end{equation}
expresses the master-to-slave latency introduced by the transceivers and
the link.
\item The recovered clock (2) is fed to a clock adjustment unit which shifts
its phase by a programmable value $phase_{S}$ to obtain the phase-compensated
clock (3) -- the final result of WR synchronization.
\item Clock (3) encodes the slave's outgoing data stream and is recovered
at the master side as signal (4).
\item The master measures the phase shift $phase_{MM}$ between its outgoing
(1) and incoming (4) clocks using a phase detector.
\end{enumerate}
The value of $phase_{MM}$ can be expressed as:
\begin{equation}
    \label{eq:phasemm}
    phase_{MM} = \{\Delta + \delta_{ms} + \delta_{sm} + phase_{S}\} \bmod
    T_{ref}
\end{equation}
where $\Delta = \Delta_{txm} + \Delta_{rxs} + \Delta_{txs} + \Delta_{rxm}$
and $T_{ref}$ is the period of Gigabit Ethernet 125 MHz reference clock (8 ns).

The goal of the presented model is to calculate the precise value of
master-to-slave offset $offset_{MS}$ by combining a coarse timestamp-based
round-trip delay \eqref{eq:meanPath} with precise phase measurement
$phase_{MM}.$ Once the offset is computed, the WR slave can phase-shift its
recovered clock (deriving $phase_{S}$ from $offset_{MS}$) to match the phase
of the master clock, completing the synchronization.
\begin{figure}[ht!]
  \centering
  \includegraphics[width=15cm]{../../../figures/protocol/sync_flow.eps}
  \caption{WR synchronization flow (DMTD is explained in Appendix~\ref{s:dmtd})}
  \label{fig:sync_flow}
\end{figure}
Determining the precise offset, however, is not a trivial task. Figure
\ref{fig:sync_flow} shows the steps needed to achieve and maintain
synchronization of a single WR link. The synchronization process can be
split into two parts: \begin{itemize}
\item \textbf{initial synchronization}, which determines the value of
$offset_{MS}$ and fixes the slave's phase shift to compensate the offset.
\item \textbf{phase tracking} which monitors the changes of $phase_{MM}$
over time and updates the phase shifter in the slave to follow these changes
and sustain synchronization.
\end{itemize}

\newpage

\subsection{Link detection and syntonization}
During the first two steps of the synchronization flow
\ref{fig:sync_flow}, a syntonized WR connection is established. Drawing
\ref{fig:link_detect_and_syntonization} shows the order of hardware operations
and message exchanges which result in a syntonized link.
\begin{figure}[ht!]
  \centering
  \includegraphics[width=7cm]{../../../figures/protocol/link_detect_and_syntonization.eps}
  \caption{WR Link detection and syntonization}
  \label{fig:link_detect_and_syntonization}
\end{figure}
At the initial moment, the master and the slave are not connected to
each other. Their PHYs are transmitting an Ethernet idle pattern (see Table 3.1 \cite{tomekMSC}), 
but are not receiving any meaningful bitstreams. As soon
as the physical connection is present, the PHYs will start receiving valid
idle patterns and the Ethernet link will become active (see \cite{IEEE802.3}
section 36.2.5.2.6). Presence of a valid physical link will trigger the
following sequence:
\begin{enumerate}
\item The master starts broadcasting \textit{ANNOUNCE} messages to look for
a WR slave,
\item Eventually, the slave will respond with a \textit{SLAVE\_PRESENT}
message, indicating that it supports WR. If no response has been received
within a predefined time, the master assumes that the slave is not WR
compatible and aborts the synchronization process,
\item The master issues a \textit{LOCK} message, commanding the slave to
begin recovering the clock from its received data stream,
\item The slave sets up its PLL to use the RX clock as a reference and as
soon as the PLL is locked, responds with a \textit{LOCKED} message.
\end{enumerate}
Upon successful completion of the above process, both nodes are syntonized
and the delay measurements can begin.

\subsection{Coarse delay measurement}
\label{s:coarse_delay}
The coarse delay measurement produces a set of timestamps $t_{1} ... t_{4}$
which are obtained using the PTP End-to-End, two-way clock method as shown
in fig. \ref{fig:wrPTPmsgs}. The timestamps are generated in the hardware
with a circuit depicted in fig. \ref{fig:coarse_measurement}, guaranteeing
single-cycle timestamping accuracy, which is necessary for reliably merging the
coarse delay with the $phase_{MM}$ phase shift during fine delay calculation.
\begin{figure}[ht!]
  \centering
  \includegraphics[width=17cm]{../../../figures/protocol/coarse_measurement.eps}
  \caption{Structure of a WR timestamping unit (TSU)}
  \label{fig:coarse_measurement}
\end{figure}
Timestamps in WR are taken when the PCS detects a Start-of-Frame Delimiter
(SFD) character in the incoming (RX timestamps) or outgoing (TX timestamps)
data stream (see section 3.2.2 \cite{tomekMSC}).

Let's first focus on the blocks marked blue in
Figure~\ref{fig:coarse_measurement}. Each time an SFD character is detected,
the PCS produces a timestamp trigger pulse which causes the timestaming unit
to take a snapshot of the free running counter CNTR\_R with the D-type
register DREG\_R. The counter is counting from 0 to 124999999 which (given
the reference clock frequency of 125 MHz) gives a period of one full second.

The counter CNTR\_R works synchronously to the reference clock (master side,
signal 1 in Figure~\ref{fig:link_model}) or the compensated clock (slave side,
signal 5). Since RX trigger pulses come from different clock domains (2 or
4), they need to be synchronized to the reference clock with a chain of
D flip-flops (SYNC). Single-cycle long trigger pulses are widened by the
pulse extender before going to the synchronizer chain to ensure that no
pulses are missed due to the metastability of synchronizer flip-flops. The
counter value is latched in register DREG\_R on the rising edge of the
synchronizer output. TX timestamp triggers, which are generated within the
reference clock domain (clocks 1 or 5) also pass through a synchronizer
chain to obtain identical trigger reaction latency.

Unfortunately, due to the jitter of clock signals, crossing clock domains
can make the gathered timestamps useless by causing a random $\pm 1$ LSB
error when the RX clock and the reference clocks are almost in phase. The
problem is illustrated in Figure~\ref{fig:ts_jitter}.  Dashed lines show
the transitions of ideal (jitter-free) clocks. If the jitter is neglected,
the reference clock transition should be slightly ahead of the RX clock
transition and the gathered timestamp should be equal to 2. However, if the
clocks are jittery, the transitions may sometimes occur in reverse order,
producing an erroneous timestamp of value 1.
\begin{figure}[ht!]
  \centering
  \includegraphics[width=10cm]{../../../figures/protocol/ts_jitter.eps}
  \caption{Timestamping errors caused by clock jitter}
  \label{fig:ts_jitter}
\end{figure}
One possible way of addressing this issue is to take RX timestamps on both
reference clock edges. The falling edge part of the TSU is marked pink in
\ref{fig:coarse_measurement}. It does not have an independent counter --
instead, the current value of the rising edge counter is latched in register
CNTR\_F on the falling edge of the reference clock, making the CNTR\us
F a copy of CNTR\_R delayed by a half of the clock period. This method
ensures that at least one of the timestamps is valid at any moment 
(see Figure~\ref{fig:ts_dualedge}). The correct timestamp is chosen depending on the
current phase shift between clocks (see section \ref{s:fine_delay}).
\begin{figure}[ht!]
  \centering
  \includegraphics[width=12cm]{../../../figures/protocol/ts_dualedge.eps}
  \caption{Dual-edge timestamping in WR}
  \label{fig:ts_dualedge}
\end{figure}
In order to simplify the hardware design, the presented TSU can only measure
the sub-second part of PTP timestamps. The UTC part is appended in the
software using the algorithm shown in listing \ref{lst:utc_gen}.
\begin{lstlisting}[caption=Producing UTC timestamps,label=lst:utc_gen]
  timestamp tstamp = get_hardware_timestamp();
  time t = get_utc_time();
// check if there was a transition of the UTC counter between the
// generation of the timestamp and its readout by the software
  if(t.milliseconds < tstamp.rising_edge)
   tstamp.seconds = t.seconds - 1;
  else
   tstamp.seconds = t.seconds;
\end{lstlisting}


\subsection{Digital DMTD phase detector}
\label{s:dmtd}
The fine delay measurements in WR are based on a Dual Mixer Time Difference
(DMTD) phase detector. Therefore, in this section a short introduction to
DMTD technology will be presented.
\begin{figure}[ht!]
  \centering
  \includegraphics[width=11cm]{../../../figures/misc/analog_dmtd.eps}
  \caption{Structure of an analog DMTD phase detector}
  \label{fig:analog_dmtd}
\end{figure}
Figure \ref{fig:analog_dmtd} shows an analog DMTD system \cite{allan90}. Let's
assume that its input clocks $a(t)$ and $b(t)$ have identical amplitudes and
frequencies equal to $f_{clk}$. The phases of both clocks are respectively
$\phi_{a}$ and $\phi_{b}$. The clocks are multiplied by the mixers with
the signal $c(t)$ of frequency $f_{offset}$ and phase $\phi_{offset}$. The
multiplication result can be expressed as:

\begin{eqnarray}
    a(t) \cdot c(t) & = & cos(2\pi t f_{clk} + \phi_{a}) \cdot cos (2\pi t
    f_{offset} + \phi_{offset}) \nonumber \\
\label{eq:mixing1}& = & \frac{1}{2} cos(2\pi t (f_{clk} + f_{offset}) +
\phi_a + \phi_{offset}) \\
\label{eq:mixing2}& + & \frac{1}{2} cos (2 \pi t (f_{clk} - f_{offset}) +
\phi_a - \phi_{offset})
\end{eqnarray}

Term \ref{eq:mixing1} has higher frequency than the multiplied signals and
it's removed by the low-pass filter, leaving only the low-frequency term
\ref{eq:mixing2}. The downconversion process changes the frequency of the
signals, but does not change their phase difference. Therefore, if the offset
frequency $f_{offset}$ is close to the input signal frequency $f_{clk}$,
the phase shift can be measured by counting the time between the rising
edges of the downconverted clocks. For example, a reference clock of 125
MHz and an offset clock of 124.99 MHz will produce an output signal of 10
kHz. At such frequencies, the phase shift can be very accurately measured
using a simple counter.

Analog DMTDs provide excellent resolution and linearity, at the cost of several
external discrete components (mixers and filters), which can be troublesome,
especially in multi-port applications such as the WR switch. Fortunately, the
analog mixing operation can be transformed into a digital sampling operation,
resulting in a digital DMTD detector, shown on fig. \ref{fig:digital_dmtd}.
\begin{figure}[ht!]
  \centering
  \includegraphics[width=12cm]{../../../figures/misc/dmtd.eps}
  \caption{Structure of a digital DMTD phase detector}
  \label{fig:digital_dmtd}
\end{figure}
In a Digital DMTD (DDMTD) \cite{icalepcs09}, the input signals are square
waves, the mixers are replaced with simple D-type flip-flops and the offset
clock is generated with a PLL from one of the input clocks. The sampling
operation performed by the flip-flops can be mathematically proved to be
equivalent to analog mixing (\ref{eq:mixing2}), but the principle of a DDMTD
can be explained in a more intuitive way.
\begin{figure}[ht!]
  \centering
  \includegraphics[width=15cm]{../../../figures/misc/dmtd_vernier.eps}
  \caption{A vernier (a) and signals generated by DDMTD (b)}
  \label{fig:dmtd_vernier}
\end{figure}
Figure \ref{fig:dmtd_vernier}a shows a vernier caliper. It has two scales -
the big millimeter scale and a small vernier scale, with the units slightly
smaller than the main scale. For a typical caliper, the vernier scale is split
into 10 intervals with 4.9 mm spacing. If the length of the measured object
has a fractional part, one of bars on the vernier scale will be in line with
one of bars from the main scale. Translating to the language of electronics,
the main scale is the input clock where each interval represents one cycle,
the vernier scale is the offset clock and the transitions in the DDMTD output
signals occur when bars on both scales are aligned. An example of signals
produced by a DDMTD with $N = 8$ is shown in fig. \ref{fig:dmtd_vernier}b. The
output phase shift is proportional to the input phase shift $\phi$ by a
factor of $(N+1)$. A general relation between the phase shift and the cycle
interval between the DMTD outputs can be expressed as:

\begin{equation}
\label{eq:dmtd1}
\phi \mbox{[ns]} = \frac{n_{cycles}}{N+1} \cdot \frac{1}{f_{in}}
\end{equation}


Digital DMTDs have all the advantages of their analog counterparts while
requiring only one external component (the oscillator producing the offset
clock) which is shared among all measurement channels.	This opens the way
for low-cost FPGA implementations. If a picosecond-level accuracy is not
required, even a PLL integrated inside an FPGA can be used, eliminating all
external components.
\begin{figure}[ht!]
  \centering
  \includegraphics[width=8cm]{../../../figures/misc/dmtd_glitches.eps}
  \caption{Glitches in the DMTD output caused by clock jitter}
  \label{fig:dmtd_glitches}
\end{figure}
In practical DDMTD implementations, the output signals need to be additionally
conditioned as the input clock jitter can introduce glitches, as shown
on fig. \ref{fig:dmtd_glitches}. More details about the deglitching and
postprocessing algorithm can be found in section 4.3.5 of \cite{tomekMSC}.

\subsection{Fine delay measurement}
\label{s:fine_delay}
Having the knowledge of PTP timestamps and round-trip phase shift $phase_{MM}$,
we can calculate the precise round-trip delay $delay_{MM}$. The calculation
is performed in two steps:
\begin{itemize}
\item The accuracy of PTP timestamps $t_{2}$ and $t_{4}$ is enhanced beyond a
single clock cycle using the knowledge of the round-trip phase $phase_{MM}$
and the slave PLL setpoint $phase_{S}$. As a result, precise timestamps
$t_{2p}$ and $t_{4p}$ are obtained.
\item The fine round-trip delay is calculated using the standard PTP formula:
\begin{equation}
\label{eq:ptp_precise}
delay_{MM} = (t_{4p} - t_{1}) - (t_{3} - t_{2p})
\end{equation}
\end{itemize}

Only the reception timestamps need to be enhanced, as they are generated
within clock domains asynchronous to the reference (or compensated) clock
(see table \ref{tab:ts_domains}). Transmission timestamps are always integer
because packets are transmitted and timestamped using the same clock.

The flow graph of the algorithm used to merge PTP timestamps with
phase measurements is shown on Figure~\ref{fig:merging_timestamps}. 
Figure~\ref{fig:merging_example} depicts sample measurements of inputs and results
of the enhancement algorithm for $t_{4p}$ timestamps, where the varying link
delay was simulated using the slave's phase shifter.

\begin{table}[htb]
  \caption{Timestamping clock domains}
  \centering
  \begin{tabular}{|c|c|c|c|c|} \hline
    Timestamp & Origin & Trigger clock & Timestamping clock & Correction \\
    \hline \hline
    $t_{1}$ & Master TX & reference (1) & reference (1) & 0 \\ \hline
    $t_{2p}$ & Slave RX & slave recovered (2) & compensated (3) & $-phase_{S}$
    \\ \hline
    $t_{3}$ & Slave TX & compensated (3) & compensated (3) & 0 \\ \hline
    $t_{4p}$ & Master RX & master recovered (4) & reference (1) & $+phase_{S}$
    \\ \hline
  \end{tabular}
  \label{tab:ts_domains}
\end{table}
\begin{figure}[ht!]
  \centering
  \includegraphics[width=13cm]{../../../figures/protocol/merging_timestamps.eps}
  \caption{Algorithm for enhancing coarse timestamps with DMTD phase.}
  \label{fig:merging_timestamps}
\end{figure}
The first step of the algorithm addresses the problem of glitches in
RX timestamps, by choosing either the rising or falling edge timestamp
depending on the phase between the RX clock and the reference clock. The
key parameter of this step, $\phi_{trans}$ provides an approximate
value of $phase_{MM}$ at which a transition should occur in the value of
$t_{4r}$. In Figure~\ref{fig:merging_example}, it is equal to 6.6 ns - it's
the approximate intersection point of $phase_{MM}$ (blue, sawtooth-like
trace) and the transition of the rising edge timestamp (red trace). The
value of $\phi_{trans}$ is a device-specific constant. It can be determined
once during the factory calibration of a WR device or measured upon every
startup by sweeping a full clock period using the built-in phase shifter
and searching for a transition in the RX timestamp value.

If the actual value of $phase_{MM}$ lies within a $\pm 25\%  T_{ref}$
range from the transition point $\phi_{trans}$ (green zone in
\ref{fig:merging_example}), the algorithm will use the $t_{4f}$ timestamp,
otherwise the $t_{4r}$ timestamp will be taken (red zone). Note that because
$phase_{MM}$ is bounded to $\left[0, T_{ref}\right)$, the range checks must
be aware of the jump in the phase between $T_{ref}$ and $0$.

The second step checks if the $phase_{MM}$ value is ahead of the transition
point $\phi_{trans}$, and eventually increases the chosen timestamp by a full
cycle ($T_{ref}$). That ensures the transition in $t_{4}$ will always occur
when $phase_{MM} = \phi_{trans}$, eliminating the risk of transition glitches.

The last part of the algorithm simply adds the picosecond part (which is
the DMTD phase corrected with the transition offset $\phi_{trans}$) to
the coarse deglitched timestamp $t_{4}$. If the resulting picosecond part
($\phi$) is negative, an additional full cycle is added to the result. The
final output of the merging algorithm is shown in \ref{fig:merging_example}
as the thick navy trace.
 \begin{figure}[ht!]
  \centering
  \includegraphics[width=12cm]{../../../figures/protocol/merging_example.eps}
  \caption{Example of $t_{4p}$ timestamp enhancing.}
  \label{fig:merging_example}
\end{figure}
The enhancement operation for the $t_{2}$ timestamp is done in a very
similar way by replacing $t_{4}$ with $t_{2}$ and $phase_{MM}$ with
$phase_{S}$. Note that changing the slave's phase shift $phase_{S}$ will
result in a change of the values of both $t_{2p}$ and $t_{4p}$ (see table
\ref{tab:ts_domains}). Increasing $phase_{S}$ by a certain $\delta$ will
cause $t_{4p}$ to also increase by $\delta$ (assuming that the link delay
stays constant). Simultaneously, the value of $t_{2p}$ will be decreased by
the same amount, as for the slave's RX timestamps, the timestamping clock is
$phase_{S}$ ahead of the trigger clock (so if the phase shift is increased,
the timestamp value goes down). Therefore, the calculated fine delay is not
affected by the changes of $phase_{S}$:

\begin{equation}
\label{eq:ptp_precise2}
delay_{MM} = (t_{4p} - t_{1}) - (t_{3} - t_{2p}) = t_{4p} \cancel{+ phase_{S}}
- t_{1} - t_{3} + t_{2p} \cancel{-phase_{S}}
\end{equation}


\subsection{Link asymmetry estimation}
\label{s:asymmetry}
In order to compute an accurate master-slave offset, one must determine the
asymmetry of the link delay. This asymmetry cannot be measured directly (as
a difference between the overall M-S and S-M delays), as it would require
the nodes to be synchronized prior to the measurement. Therefore it is
only possible to estimate the asymmetry from round-trip delay $delay_{MM}$,
using the knowledge of the properties of the components constituting the link.

In the WR optical link model, the following sources of asymmetry were taken
into consideration:
\begin {itemize}
\item Propagation delays of electronic components and PCB traces (circuit
asymmetry),
\item Asymmetry of optical transceivers (SFPs),
\item Difference between TX and RX wavelengths in the fiber,
\item Internal structure and clocking of the PHY (SerDes) chips.
\end{itemize}
\begin{figure}[ht!]
  \centering
  \includegraphics[width=\textwidth]{../../../figures/protocol/asymmetries.eps}
  \caption{Delay asymmetries in WR optical link.}
  \label{fig:asymmetries}
\end{figure}
Figure \ref{fig:asymmetries} shows the reference asymmetry model used in WR
devices. The device's asymmetric delays (as shown in Figure~\ref{fig:link_model})
are expressed as sums of circuit, SFP and PHY delays between the phase detector
inputs (the PD measuring $delta_{MM}$ on the master side and the PD in the
phase shifting PLL on the slave side) and the SFP optical input/output. The
fiber asymmetry is compensated separately in the slave's PTP servo.
\begin{eqnarray}
\Delta_{tx(m/s)} = \delta_{TX\us PHY(M/S)} + \delta_{TX\us CIR(M/S)} +
\delta_{TX\us SFP(M/S)}
\\
\Delta_{rx(m/s)} = \delta_{RX\us PHY(M/S)} + \delta_{RX\us CIR(M/S)} +
\delta_{RX\us SFP(M/S)} \nonumber
\label{eq:asym_delay}
\end{eqnarray}\\

\subsubsection{Circuit and SFP asymmetry}
Circuit asymmetry stems from differing PCB trace lengths, propagation delays
of the electronic components of the clock distribution system and FPGA logic
cell placement and routing delays. SFP asymmetry is a result of different
reception and transmit delays between the electrical and optical ports of
the SFP transceiver.

These asymmetries can vary with changes of the device's operating conditions
(i.e. temperature, supply voltage). Depending on the synchronization accuracy
requirements, they can be:
\begin{itemize}
\item Treated as time-invariant and measured once during factory setup.
\item Actively compensated, following the changes in temperature and/or supply
voltage read from built-in sensors (requires a model of delay vs. temperature
and voltage),
\item Eliminated by building a system where both master and slave have
identical (or zero) asymmetry and operate in similar conditions.
\end{itemize}

The last method is particularly useful for compensating the delays inside
the FPGA on the way between the PHY (oscillator) and the phase detector,
as they cannot be externally measured during the factory calibration
process. They are also significantly affected by temperature and voltage
changes. Therefore, the only practical way of dealing with the FPGA part
of the circuit asymmetry is to equalize the delays on all phase detector
inputs. This can be done by constraining the routing delays or manually
placing and routing the phase detector block. This approach also allows for
reducing the impact of the temperature-voltage changes on the asymmetry to
a negligible value, as two equal paths placed next to each other would have
very similar temperature/voltage delay coefficients.

SFP asymmetries can be compensated in a similar manner by choosing pairs
of SFP transceivers whose TX and RX delays differ by similar values
(i.e. $\delta_{TX\us SFPM} - \delta_{RX\us SFPM} = \delta_{TX\us SFPS} -
\delta_{RX\us SFPS}$). This can be done in a laboratory system where all
other asymmetric delays are already compensated by selecting SFPs that
together give minimum clock offset.

\subsubsection{WDM fiber asymmetry}
\label{s:wdm_asymmetry}
As mentioned in section 3.2 of \cite{tomekMSC}, WR uses different wavelengths
for transmitting and receiving the data (for example 1550/1310 nm). Due
to chromatic dispersion of the fiber, the refractive indexes for these
wavelengths are slightly different, causing different propagation speeds
(and thus, different delays) between the master and the slave.

The refractive index at a given wavelength can be derived using Sellmeier's
equation  \cite{saleh07} (\ref{eq:sellmeier}):

\begin{equation}
\label{eq:sellmeier}
n^2(\lambda) = 1 + \frac{B_{1} \lambda^2}{\lambda^2 - C_{1}} + \frac{B_{2}
\lambda^2}{\lambda^2 - C_{2}} + \frac{B_{3} \lambda^2}{\lambda^2 - C_{3}}
\end{equation}

where $B_{i}$ and $C_{i}$ are material-specific coefficients.

For a standard G.652 telecom fiber, the refractive indexes are respectively:
$n_{1550} = 1.467$ and $n_{1310} = 1.466$. In order to simplify the asymmetry
calculations in the hardware, the WR specification defines
a custom fiber asymmetry coefficient \eqref{eq:fiber_alpha1} expressing the
ratio between the M-S and S-M fiber propagation delays:

\begin{equation}
\label{eq:fiber_alpha1}
\alpha = \frac{\delta_{ms}}{\delta_{sm}} - 1 = \frac{n_{1550}}{n_{1310}} - 1
\end{equation}
Unfortunately, refractive indexes may vary slightly between different fiber
manufacturers, making the direct calculation of $\alpha$ unreliable. WR
can solve this problem by characterizing the fiber asymmetry using
laboratory measurements of $delay_{MM}$ (done by PTP) and clock offset
$offset_{MS}$ (done using an oscilloscope). The measurements are performed
when all the other asymmetries are already compensated. The value of
$\alpha$ \ref{eq:fiber_alpha3} is calculated by solving equation system
\ref{eq:fiber_alpha2}:
\begin{equation}
\label{eq:fiber_alpha2}
\begin{cases}
delay_{MM} = \delta_{ms} + \delta_{sm} + \Delta \\
offset_{MS} = \frac{\delta_{ms} - \delta_{sm}}{2}
\end{cases}
\end{equation}
where $\Delta$ accounts for all fixed delays in the path (i.e. $\Delta_{txm}
+ \Delta_{rxs} + \Delta_{txs} + \Delta_{rxm}$). Substitution of $\delta_{ms}$
and $\delta_{sm}$ with the results of \ref{eq:fiber_alpha2} gives the final
value of $\alpha$:
\begin{equation}
\label{eq:fiber_alpha3}
\alpha = \frac{delay_{MM}  - \Delta + 2\cdot offset_{MS}}{delay_{MM} -
\Delta - 2\cdot offset_{MS}}
\end{equation}

\subsubsection{Transceiver asymmetry}
\label{s:xcvr_asymmetry}
\label{sec:calibForGigbitE}
\textit{[This chapter has been modified to add clarity]}


Transceiver asymmetry is a result of the internal structure of the SerDes's
circuitry. Most of the SerDes chips nowadays are optimized for low power
consumption and fast locking to the incoming data stream. Unfortunately,
because of these optimizations, PHYs may not keep constant transmit/receive
latencies. The problem is illustrated in Figure~\ref{fig:phy_asymmetry}.
\begin{figure}[ht!]
  \centering
  \includegraphics[width=11cm]{../../../figures/protocol/phy_asymmetry.eps}
  \caption{Random delays in gigabit SerDes devices (a) and blocks causing them
  (b).}
  \label{fig:phy_asymmetry}
\end{figure}
PHY asymmetry manifests itself as a random latency between the rising edge of
the TX/RX clock and the inter-symbol boundary in the transmitted (received)
serial data stream. In most PHYs, this latency can be different for each
PLL/CDR lock cycle, but once the PHY is locked, the delays shall remain
constant. PHYs whose TX/RX delays can change when the link is active are
unsuitable for WR devices.

In the PHYs which have been evaluated for usability in WR hardware (TLK1221
from Texas Instruments and Virtex-6/Spartan-6 GTP transceivers), the following
random delays were identified:
\begin{itemize}
\item RX alignment latency, observed in PHYs which do not correct the RX
clock phase when aligning to the inter-symbol boundary (red, observed for
Xilinx GTP),
\item RX latency resulting from the internal structure of the digital
oversampling CDR (green, observed for TLK1221),
\item TX latency caused by the clock divider between the internal PLL and
the parallel-to-serial register (blue, observed for TLK1221).
\end{itemize}
Alignment latency can be measured every time the link goes up by disabling
automatic comma alignment and bit-shifting the unaligned output data
until a valid 8B10B code sequence is detected (\textit{bitslip trick},
see \cite{Peek2010}). Compensation of the latter two latencies, however,
requires additional calibration logic. An example method (used in the WR
switch) is shown in Figure~\ref{fig:phy_latency_measurement}.
\begin{figure}[ht!]
  \centering
%  \includegraphics[width=\textwidth]{fig/tomeksDrawings/phy_latency_measurement.eps}
  \includegraphics[width=0.60\textwidth]{../../../figures/misc/calibration_1.ps}
  \caption{PHY latency measurement using calibration patterns.}
  \label{fig:phy_latency_measurement}
\end{figure}
Such methods can by applied to the
PHYs whose maximum ($\delta_{PHY(M/S)_{max}}$) and minimum ($\delta_{PHY(M/S)_{min}}$) delays are known 
(provided in the datasheet) and the delay's variation is
\begin{equation}
  \label{eq:fixedDelayVariation}
  \delta_{\{TX, RX\}\_PHY(M/S)_{variable}} \in \langle 0 :\delta_{PHY(M/S)_{max}} - \delta_{PHY(M/S)_{min}} \rangle
\end{equation}
e.g. below 10 bit times in the case of Gigabit Ethernet. 
Therefore, a fixed delay can be expressed as a sum of a constant value 
($\delta_{\{TX, RX\}\_PHY(M/S)_{min}}$) and a variable part ($\delta_{\{TX, RX\}\_PHY(M/S)_{variable}}$) which needs to be measured:
\begin{equation}
  \label{eq:fixedDelay}
  \delta_{\{TX, RX\}\_PHY(M/S)} = \delta_{\{TX, RX\}\_PHY(M/S)_{min}} + \delta_{\{TX, RX\}\_PHY(M/S)_{variable}}
\end{equation}
 


The PHY transmit path is fed with a sequence of RD+ K28.5 characters (1111100000),
effectively producing a 125 MHz square wave on the serial outputs. 
The repeated pattern of five "0" and five "1" is defined by the IEEE 802.3 standard \cite{IEEE802.3} 
as \textit{Low-frequency test pattern} (Appendix~36A.2). The phase
shift between the TX clock and the output bitstream can be measured using
a DDMTD phase detector, giving the value of the TX variable latency 
($\delta_{TX\_PHY(M/S)_{variable}}$). The same method
can be used to measure the RX variable latency ($\delta_{RX\_PHY(M/S)_{variable}}$) 
by commanding the link partner to send
the calibration pattern. Note that since the K28.5 character contains a comma
(5 consecutive ones), a burst of subsequent K28.5 symbols will cause improper
operation of the PHY's comma alignment unit. Therefore, comma alignment must
be disabled during the calibration process.

\subsection{Establishing and maintaining synchronization}
\label{s:estab_sync}
Having obtained the values of round-trip delay $delay_{MM}$ and link asymmetry,
we can calculate the one-way master to slave delay $delay_{MS}$ by solving
equation (\ref{eq:delaymm_full}):

\begin{equation}
delay_{MM} = \Delta + \delta_{ms} + \delta_{sm}
\label{eq:delaymm_full}
\end{equation}

where $\Delta = \Delta_{txm} + \Delta_{rxs} + \Delta_{txs} +
\Delta_{rxm}$. Substituting $\delta_{sm}$ with \ref{eq:fiber_alpha1}
and solving for $\delta_{ms}$, we obtain the one-way fiber delay
\ref{eq:delaymm_full_2}:

\begin{equation}
\delta_{ms} = \frac{1+\alpha}{2+\alpha} (delay_{MM} - \Delta)
\label{eq:delaymm_full_2}
\end{equation}

which after adding the circuit, SFP and PHY asymmetric delays present on the
master to slave path gives the final master-slave delay \ref{eq:delaymm_full_3}
and offset \ref{eq:offset_ms}:

\begin{eqnarray}
\label{eq:delaymm_full_3}delay_{MS} & = & \frac{1+\alpha}{2+\alpha} (delay_{MM}
- \Delta) + \Delta_{txm} + \Delta_{rxs} \\
\label{eq:offset_ms} offset_{MS} & = & t_{1} - t_{2p} - delay_{MS}
\end{eqnarray}
The value of $offset_{MS}$ is the input for the slave's offset adjustment
algorithm which controls the slave's clock servo. The flow diagram of the
algorithm is shown in fig. \ref{fig:adjustment_and_servo}a and an example
servo design can be found in fig. \ref{fig:adjustment_and_servo}b. The
algorithm assumes that the frequency has been already syntonized by means
of Sync-E and only the clock offset needs to be corrected. Offset correction
is split into 3 steps:
\begin{figure}[ht!]
  \centering
  \includegraphics[width=\textwidth]{../../../figures/protocol/adjustment_and_servo.eps}
  \caption{WR slave offset adjustment (a) and clock servo (b)}
  \label{fig:adjustment_and_servo}
\end{figure}
\begin{enumerate}
\item \textbf{UTC time adjustment}: the UTC counter in the servo is
increased (or decreased) by the number of full seconds \ref{eq:corr_utc}
in $offset_{MS}$:
\begin{equation}
corr_{utc} = \lfloor{\frac{offset_{MS}}{1 \mathrm{~s}}}\rfloor
\label{eq:corr_utc}
\end{equation}
\item \textbf{Clock cycle counter adjustment}: The reference clock cycle
counter, which produces the PPS signal is adjusted by the number of full
$T_{ref}$ (8 ns) cycles \ref{eq:corr_cnt}.
\begin{equation}
corr_{cnt} = \lfloor \frac{offset_{MS} - corr_{UTC}}{T_{ref}}\rfloor
\label{eq:corr_cnt}
\end{equation}
\item \textbf{Phase adjustment}: The slave's phase shifter is adjusted with
the remaining sub-cycle part of the offset \ref{eq:corr_phase}
\begin{equation}
corr_{phase} = offset_{MS} - [offset_{MS}]
\label{eq:corr_phase}
\end{equation}
\end{enumerate}
\emph{Voil\`{a}!} Now the slave's clock and PPS signals are synchronized to the
master with sub-nanosecond accuracy. Since the offset can vary with operating
conditions, it is measured at regular intervals and the difference between
subsequent measurements is added to the slave's phase shift to compensate for
phase drift:
\begin{equation}
corr_{phase} = offset_{MS} - offset_{MS\us previous}
\label{eq:corr_phase2}
\end{equation}
Because the phase drift is mainly caused by temperature variations, the rate
of subsequent adjustments can be very low, even in the scale of a single
adjustment per hour.

Note that the way $corr_{phase}$ is calculated requires the phase shifter
to be able to change the phase relatively by any arbitrary value with
no "jumps" in the signal when the value of $corr_{phase}$ crosses the
inter-cycle boundary. For example, one can use a PLL with a phase detector
capable of handling wrap-around phase transitions, but not a programmable
delay line. An example design of such phase detector and PLL is described
in section 4.3.7 of \cite{tomekMSC}.


%%%%%%%%%%%%%%%%%%%%%%%%%%%%%%%%%%%%%%%%%%

% \subsection{Measurement of fixed delays for Gigabit Ethernet over Optic Fiber}
% \label{sec:calibForGigbitE}
% 
% The variation of $\Delta_{\{tx_m, rx_s, tx_s, rx_m\}}$ delays is often caused by 
% the PHY's serializer / deserializer (SerDes), phase locked loop (PLL) or clock and 
% data recovery circuitry (CDR). The delay on the PHY can be measured by detecting 
% the phase shift between SerDes I/O and Tx/Rx clock. Such methods can by applied to the
% PHYs whose maximum ($\Delta_{max}$) and minimum ($\Delta_{min}$) delays are known 
% (provided in the datasheet) and the delay's variation is
% \begin{equation}
%   \label{eq:fixedDelayVariation}
%   \Delta_{variable} \in \langle 0 : \Delta_{max} - \Delta_{min} \rangle
% \end{equation}
% e.g. below 10 bit times in case of Gigabit Ethernet. 
% Therefore, a fixed delay can be expressed as a sum of a constant value 
% ($\Delta_{min}$) and a variable part ($\Delta_{variable}$) which needs to be measured:
% \begin{equation}
%   \label{eq:fixedDelay}
%   \Delta_{\{tx_m, rx_s, tx_s, rx_m\}} = \Delta_{min\_\{tx_m, rx_s, tx_s, rx_m\}} + \Delta_{variable\_\{tx_m, rx_s, tx_s, rx_m\}}
% \end{equation}
%  
% In White Rabbit, the measurement of fixed delay's variable part ($\Delta_{variable}$) is done by 
% sending a repeated pattern of five "0" and five "1" (0000011111) over Gigabit Ethernet.
% Such signal creates a 125~MHz clock on the SerDes I/O. Since the Tx/Rx clock 
% frequency is 125~MHz, the phase shift between the SerDes I/O and the Tx/Rx 
% clocks is equal to $\Delta_{variable}$ of the PHY (see Figure~\ref{fig:wrCalibration}). 
% 
% 
% The repeated pattern of five "0" and five "1" is defined by the IEEE 802.3 standard \cite{IEEE802.3} 
% as \textit{Low-frequency test pattern} (Appendix~36A.2). It shall be generated in WR nodes 
% by a repeated transmission of RD+ K28.7 code-group.
% 
% 
% \begin{figure}[ht!]
%   \centering
%   \includegraphics[width=0.60\textwidth]{fig/calibrate.ps}
%   \caption{Measurement of fixed delays $\Delta_{\{tx, rx\}}$ in Gigabit Ethernet-based WR node with 
% not fully-deterministic PHY.}
%   \label{fig:wrCalibration}
% \end{figure}
% 
% Measurement of fixed delays for Gigabit Ethernet over optic fiber, and any other medium, is optional. 
% It is not needed if deterministic PHYs or internal FPGA transceivers which can be internally 
% characterized \cite{Peek2010} are used. In such case, the information about the fixed delays is 
% distributed across the link without preceding measurement.
% 
% \newpage
% 
% \subsection{Fine Delay Measurement}
% \label{sec:FineDelayMeasurement}


\newpage

\section{Typical flow of WR Signaling message exchange}
\label{ap:wr_lsu_flow}
\begin{figure}[ht!]
  \centering
%  \vspace{-1.3cm}
  \includegraphics[width=1.0\textwidth]{protocol/wrMSGsExchangeFlow.ps}
  \caption{Typical flow of events (no exceptions) during WR Link Setup.}
  \label{fig:wrFSMcommun}
\end{figure}

\newpage


\section{PTP and WR FSMs from POWER ON use case}
\label{ap:ptpAndWrFSMS}
\begin{figure}[ht!]
  \centering
%  \vspace{-1.3cm}
  \includegraphics[width=1.0\textwidth]{protocol/ptpAndWrFSMs.ps}
  \caption{PTP and WR FSMs from POWER ON use case}
  \label{fig:wrFSMcommun}
\end{figure}

\newpage


\section{Primitive data types}
\label{ap:wrDataTypes}

\begin{table}[ht!]
\caption{Primitive data types}
\centering
\begin{tabular}{| c | c |}          					\hline
\textbf{Data type  } &  \textbf{Definition}				  \\ 
		     &  	               				  \\ \hline
UInteger8	     &   8-bit unsigned integer				  \\ \hline
UInteger16	     &  16-bit unsigned integer				  \\ \hline
UInteger32	     &  32-bit unsigned integer				  \\ \hline
 Integer64	     &  64-bit integer					  \\ \hline
UInteger64	     &  64-bit unsigned integer				  \\ \hline
\end{tabular}
\label{tab:wrPtpProfile}
\end{table}


\section{Computations using the two-step delay request-response mechanism with asymmetry correction}
\label{ap:computationsExplanation}

The method of incorporating asymmetry obtained using the WR Link Model by the WR Slave 
is WRPTP-implementation specific. Below, a summary of PTP offset ($<offsetFromMaster>$) 
and mean delay ($<meanDelayPath>$) calculations are presented. They are followed by 
two methods of including WR asymmetry, with the later being recommended. 
The method used might vary among WR Slaves.

\subsection{Overview of PTP offset and mean path delay calculations}
\label{ap:ptpComputations}
[Citing \cite{IEEE1588}] \\
The two-step delay request-response for a boundary or ordinary clock,:
\begin{enumerate}
  \item Sync message transmission (section 9.5.9 of \cite{IEEE1588}):
    \begin{itemize}
      \item The $<syncEventEgressTimestamp>$  ($t_1$) shall be generated upon 
	    transmission of the Sync message.
      \item The originTimestamp field of the Sync message shall be 0.
      \item The correctionField of the Sync message shall be set to 0.
    \end{itemize}
%%%%%%%%%%%%%%%%%%%%%%%%%%%%%%%%%%%%%%%%%%%%%%%%%%%%%%%%%
	\begin{align}
	  \label{eq:Sync_tx}
	  Sync.originTimestamp = 0 \\
	  Sync.correctionField = 0
	\end{align}
%%%%%%%%%%%%%%%%%%%%%%%%%%%%%%%%%%%%%%%%%%%%%%%%%%%%%%%%%
  \item Sync message reception (section 9.5.4 and 11.6.2 of \cite{IEEE1588}):
    \begin{itemize}
      \item The $t_2$ shall be generated upon reception of the Sync message.
      \item The sync message correctionField shall be adjusted for asymmetry by adding the value of 
	    the ingress path delayAsymmetry prior to its any use in a computation.
    \end{itemize}
%%%%%%%%%%%%%%%%%%%%%%%%%%%%%%%%%%%%%%%%%%%%%%%%%%%%%%%%%
	\begin{align}
	  \label{eq:Sync_rx}
	  Sync.correctionField = delayAsymmetry
	\end{align}
%%%%%%%%%%%%%%%%%%%%%%%%%%%%%%%%%%%%%%%%%%%%%%%%%%%%%%%%%
  \item Follow\_Up message transmission (section 9.5.10 of \cite{IEEE1588}):
    \begin{itemize}
      \item The preciseOriginTimestamp field of the Follow\_Up should be \\
	    the $<syncEventEgressTimestamp>$ of the associated Sync message excluding 
	    any fractional nanoseconds of the associated Sync message.
      \item The correctionField field of the Follow\_Up shall be the fractional nanoseconds part of 
	    the $<syncEventEgressTimestamp>$ of the associated Sync message.
    \end{itemize}
%%%%%%%%%%%%%%%%%%%%%%%%%%%%%%%%%%%%%%%%%%%%%%%%%%%%%%%%%
	\begin{align}
	  \label{eq:follow_up_tx}
	  Follow\_Up.preciseOriginTimestamp &= seconds\_and\_nanoseconds(t_1)\\
	  Follow\_Up.correctionField &=  fractional\_ns\_part(t_1)\\
	  <syncEventEgressTimestamp>= t_1 &= Follow\_Up.preciseOriginTimestamp \\
		\nonumber	     &+ Follow\_Up.correctionField
	\end{align}
%%%%%%%%%%%%%%%%%%%%%%%%%%%%%%%%%%%%%%%%%%%%%%%%%%%%%%%%%
  \item Delay\_Req message transmission (section 11.6.3 and 11.3.2 of \cite{IEEE1588}):
    \begin{itemize}
      \item Prior to transmission on an egress port the correctionField of 
	    the Delay\_Req message shall be modified by subtracting the value of the egress 
	    path delayAsymmetry from its correctionField.
      \item The $t_3$ timestamp shall be generated upon transmission of the Delay\_Req message.
    \end{itemize}
%%%%%%%%%%%%%%%%%%%%%%%%%%%%%%%%%%%%%%%%%%%%%%%%%%%%%%%%%
	  \begin{align}
	    \label{eq:follow_up_rx}
	    Delay\_Req.correctionField &= -delayAsymmetry \\
	    Delay\_Req.originTimestamp &= 0
	  \end{align}
%%%%%%%%%%%%%%%%%%%%%%%%%%%%%%%%%%%%%%%%%%%%%%%%%%%%%%%%%
  \item Delay\_Req message reception (section 9.5.6 of \cite{IEEE1588}):
    \begin{itemize}
      \item The $<delayReqEventIngressTimestamp>$ ($t_4$)  
	    shall be generated upon receipt of the Delay\_Req message.
    \end{itemize}
  \item Delay\_Resp message transmission (section 9.5.12 and 11.3.2 of \cite{IEEE1588}):
    \begin{itemize}
      \item The receiveTimestamp field of the Delay\_Resp message shall be set to the seconds and 
	    nanoseconds portion of the time $t_4$.
      \item The correctionField of the Delay\_Req message shall be set to 0 and then 
	    the correctionField of the Delay\_Resp message shall be added to it and
	    the fractional nanoseconds portion of $t_4$ shall be subtracted from the 
	    corrctionField of Delay\_Resp message.
    \end{itemize}
%%%%%%%%%%%%%%%%%%%%%%%%%%%%%%%%%%%%%%%%%%%%%%%%%%%%%%%%%
	  \begin{align}
	    \label{eq:delay_resp_tx}
	    Delay\_Resp.correctionField 	&= Delay\_Req.correctionField \\
					  &- fractional\_ns\_part(t_4)\\
	  \nonumber	 		&= -(delayAsymmetry + fractional\_ns\_part(t_4)) \\
	    Delay\_Resp.receiveTimestamp 	&= seconds\_and\_nanoseconds(t_4)
	  \end{align}
%%%%%%%%%%%%%%%%%%%%%%%%%%%%%%%%%%%%%%%%%%%%%%%%%%%%%%%%%
  \item Delay\_Resp message reception (section 11.3.2 of \cite{IEEE1588}):
    \begin{itemize}
      \item The $<meanPathDelay>$ shall be computed as follows
    \end{itemize}
%%%%%%%%%%%%%%%%%%%%%%%%%%%%%%%%%%%%%%%%%%%%%%%%%%%%%%%%%
	  \begin{align}
	    \label{eq:meanPathDelay}
	    <meanPathDelay> &=[(t_2 - t_3) \\
	  \nonumber	  &+ (Delay\_Resp.receiveTimestamp \\
	  \nonumber	  &- Follow\_Up.preciseOeriginTrimestamp)\\ 
	  \nonumber	  &- Sync.correctionField \\
	  \nonumber	  &- Follow\_Up.correctionField \\
	  \nonumber	  &- Delay\_Resp.correctionField]/2
	  \end{align}
%%%%%%%%%%%%%%%%%%%%%%%%%%%%%%%%%%%%%%%%%%%%%%%%%%%%%%%%%
  \item Computation of the clock offset  (section 11.2 of \cite{IEEE1588}):
%%%%%%%%%%%%%%%%%%%%%%%%%%%%%%%%%%%%%%%%%%%%%%%%%%%%%%%%%
	  \begin{align}
	    \label{eq:offset}
	    <offsetFromMaster> &=t_2 \\
	  \nonumber	  &- Follow\_Up.preciseOriginTimestamp \\
	  \nonumber	  &- <meanPathDelay> \\
	  \nonumber	  &- Sync.correctionField \\
	  \nonumber	  &- Follow\_Up.correctionField 
	  \end{align}
%%%%%%%%%%%%%%%%%%%%%%%%%%%%%%%%%%%%%%%%%%%%%%%%%%%%%%%%%
\end{enumerate}

\subsection{WR asymmetry as PTP communication path asymmetry}

The asymmetry obtained using the WR Link Model ($asymmetry$, section \ref{sec:delayAsymCal}) 
can be seen as a communication path asymmetry $delayAsymmetry$ as defined in 7.4.2 
of \cite{IEEE1588}
\begin{equation}
\label{eq:ptpAsym}
  delayAsymmetry= asymmetry
\end{equation}
and incorporated into the PTP calculations as described in Appendix~\ref{ap:ptpComputations}, 
equations \eqref{eq:Sync_rx} and \eqref{eq:follow_up_rx}.

\subsubsection{Solution for Ethernet over a Single-mode Optical Fiber}

In the case of the solution for Ethernet over a single-mode Optic Fiber 
(section\ref{sec:EthernetOverSingleModeFiber}), it can be calculated using directly 
equation~\eqref{eq:aqasymm}, where $\mu$ is $<meanDelayPath>$ obtained using 
equation~\eqref{eq:meanPathDelay}:
\begin{equation}
\label{eq:delayAsymmetry}
  delayAsymmetry= \Delta_{tx_m} + \Delta_{rx_s} - \frac{\Delta - \alpha * <meanDelayPath> 
			  + \alpha * \Delta}{2 + \alpha}
\end{equation}
It needs to be noted, that the $delayAsymmetry$ used in calculation of $<meanDelayPath>$ and\\
$<offsetFromMaster>$ is computed based on a previous measurement of the \\$<meanDelayPath>$. 
This is because the standard \cite{IEEE1588} seems to assume that, if available, the 
$delayAsymmetry$ is constant and known in advance.

Therefore, a direct incorporation of the asymmetry calculated using the WR Link Model into 
$<meanDelayPath>$ and $<offsetFromMaster>$ is recommended.

\subsection{Direct WR asymmetry incorporation into PTP computations}

The asymmetry obtained using the WR Link Model ($asymmetry$, section \ref{sec:delayAsymCal})  
can be incorporated directly into the computation 
of $<meanDelayPath>$ and $<offsetFromMaster>$. In such case, the $delayAsymmetry$ as defined 
in 7.4.2 of \cite{IEEE1588}, is set to 0x0 (equations \eqref{eq:Sync_rx} and 
\eqref{eq:follow_up_rx}). The $<meanDelayPath>$ value is calculated using \eqref{eq:meanPathDelay} 
and corrected with $asymmetry$ to obtain $<pathDelay_{MS}>$,
\begin{equation}
\label{eq:pathDelayMS}
  <pathDelay_{MS}>= <meanDelayPath> + asymmetry
\end{equation}
which is used to compute $<offsetFromMaster>$ using \eqref{eq:offset}:
\begin{align}
  \label{eq:WRoffset}
  <offsetFromMaster> &=t_2 \\
\nonumber	  &- Follow\_Up.preciseOriginTimestamp \\
\nonumber	  &- <pathDelay_{MS}> \\
\nonumber	  &- 0 \\
\nonumber	  &- Follow\_Up.correctionField \\
		  &=t_{2}-t_1 - <pathDelay_{MS}>
\end{align}


\subsubsection{Solution for Ethernet over a Single-mode Optical Fiber}

A WR Slave, which assumed communication path asymmetry to be 0 ($delayAsymmetry=0$) and measured
timestamps $t_2$ and $t_3$, can assemble the timestamps received from the WR Master follows:
\begin{align}
  \label{eq:t1}
	    t_1 &= Follow\_Up.preciseOriginTimestamp \\
 \nonumber 	&+ Follow\_Up.correctionField 
\end{align}
and
\begin{align}
  \label{eq:t4}
	    t_{4} 	&= Delay\_Resp.receiveTimestamp \\
\nonumber		&- Delay\_Resp.correctionField 	
\end{align}
Having all timestamps ($t_1$, $t_2$, $t_3$, $t_4$), the fixed delays 
($\Delta_{txm}$,$\Delta_{rxs}$, $\Delta_{txs}$, $\Delta_{rxm}$) and relative delay coefficient 
($\alpha$), the WR Link Model can be combined with the PTP computation to obtain directly the 
correct values of $<delayPath_{MS}>$ and $<offsetFromMaster>$ :
\begin{eqnarray}
  <delayPath_{MS}> & = & \frac{1+\alpha}{2+\alpha} 
			[(t_4-t_1) + (t_3-t_2) - \Delta] + \Delta_{txm} + \Delta_{rxs} \\
  <offsetFromMaster> & = & t_{2} - t_{1}  - <delayPath_{MS}>
\end{eqnarray}
where
\begin{eqnarray}
  \Delta & = & \Delta_{txm} + \Delta_{rxs} + \Delta_{txs} + \Delta_{rxm}
\end{eqnarray}

\newpage


\begin{thebibliography}{9}



\bibitem{IEEE1588}
  IEEE Std 1588-2008
  \emph{IEEE Standard for a Precision Clock Synchronization Protocol for Networked Measurement and 
Control Systems}.
  IEEE Instrumentation and Measurement Society, New York,
  2008,
  http://ieee1588.nist.gov/.

\bibitem{IEEE802.3}
  IEEE Std 802.3-2008
  \emph{IEEE Standard for Information technology -- Telecommunications and information exchange 
	between systems -- Local and Metropolitan area networks -- Specific requirements}.
  IEEE Computer Society, New York,
  2008.

\bibitem{SynchE}
  ITU-T G.8262/Y.1362
  \emph{Timing characteristics of a synchronous
  Ethernet equipment slave clock}.
  TELECOMMUNICATION STANDARDIZATION SECTOR OF ITU, 
  07/2010.

\bibitem{tomekMSC}
  Tomasz W\l{}ostowski
  \emph{Precise time and frequency transfer in a White Rabbit network}.
  Warsaw University of Technology, 
  05/2011.

\bibitem{Peek2010}
  P.P.M. Jansweijer, H.Z. Peek,
  \emph{Measuring propagation delay over a 1.25 Gbps bidirectional data link}.
  National Institute for Subatomic Physics, Amsterdam,
  2010,\\
  http://www.nikhef.nl/pub/services/biblio/technicalreports/ETR2010-01.pdf.

\bibitem{Takahide}
  Takahide Murakami and Yukio Horiuchi,
  \emph{A Master Redundancy Technique in IEEE 1588 Synchronization with a Link Congestion
  Estimation}.
  ISPCS Proccedings,
  2010,\\

\bibitem{allan90}
  D.B Sullivan, D.W. Allan, D.A. Howe, F.L. Walls,
  \emph{Characterization of Clocks and Oscillators}.
  http://tf.nist.gov/general/pdf/868.pdf, 
  NIST Technical Note,
  1990.\\

\bibitem{icalepcs09}
  J. Serrano, P. Alvarez, M. Cattin, E. G. Cota, J. H. Lewis, P. Moreira, T. W\l{}ostowski
  and others,
  \emph{The White Rabbit Project}.
  ICALEPCS TUC004,
  2009.\\

\bibitem{saleh07}
  Bahaa E.A. Saleh and Malvin Carl Teich,
  \emph{Fundamentals of Photonics}.
  Wiley-Interscience,
  2007.\\

\end{thebibliography}


\end{document}
